
\documentclass[foods,review,accept,pdftex,moreauthors]{Definitions/mdpi} 


\usepackage{lineno}
\let\linenumbers\relax
\let\internallinenumbers\relax
\nolinenumbers

%=================================================================
% MDPI internal commands - do not modify
\firstpage{1} 
\makeatletter 
\setcounter{page}{\@firstpage} 
\makeatother
\pubvolume{1}
\issuenum{1}
\articlenumber{0}
\pubyear{2026}
\copyrightyear{2026}
\externaleditor{Ali Khoddami} 
\datereceived{12 May 2026} 
\daterevised{30 May 2026} 
\dateaccepted{9 June 2026} 
\datepublished{ } 

%=================================================================
\newcommand{\todo}[1]{\textcolor{red}{[TODO: #1]}}
\usepackage[normalem]{ulem}
\usetikzlibrary{shapes.geometric, arrows.meta, positioning}

%=================================================================

\Title{\highlighting{Immature} %MDPI: Important: 1. We have done the initial layout for your manuscript through our layout team. Please do not change the layout, otherwise we cannot proceed to the next step. 2. Please note that after the article has been accepted, changes to authorship are not permitted. 3. Please do not delete our comments. 4. Please revise and answer all the comments we left. E.g.: “It should be italic”; “I confirm”; “I have checked and revised all.” 5. Please correct directly on this version. 6. The paper has been edited by our in-house English editor, please check carefully throughout the manuscript. 7. Please check the whole paper very carefully, any changes are not allowed after proofreading.
 Honey as a Quality Challenge in Global \linebreak  Apicultural Production}


\newcommand{\orcidauthorA}{0009-0004-1158-0693}
\newcommand{\orcidauthorB}{0009-0003-7214-4097} 
\newcommand{\orcidauthorC}{0000-0001-9178-1540} 
\newcommand{\orcidauthorD}{0000-0003-2580-7954} 
\newcommand{\orcidauthorE}{0000-0003-3900-7368} 
\newcommand{\orcidauthorF}{0000-0002-0754-0806}
\Author{\hl{Anna Gajda} %MDPI: 1. Please carefully check the accuracy of names and affiliations. 2. Please note that changing authorship is NOT acceptable during this period (including adding new authors/corresponding authors/affiliations or deleting present authors/corresponding authors/affiliations or exchanging author orders). 
%I confirm
 $^{1}$\orcidE{}, 
 Bartosz Lewandowski $^{2}$\orcidA{}, 
 Przemysław Rujna $^{3}$\orcidB{},
 Joanna Katarzyna Banach  $^{4,}$*\orcidC{},
 \linebreak  Renata Pietrzak-Fiećko $^{5}$\orcidF{} and
 Ewaryst Tkacz \hl{$^{2}$}%MDPI: We merged the duplicated address into one affiliation, please confirm.
\orcidD{}}


\AuthorNames{Anna Gajda, Bartosz Lewandowski, Przemysław Rujna, Joanna Katarzyna Banach,  Renata Pietrzak-Fiećko, Ewaryst Tkacz}


\address{%
$^{1}$ \quad Laboratory of Bee Diseases, Institute of Veterinary Medicine, Warsaw University of Life Sciences, \linebreak  02-787 Warsaw, Poland; anna\_gajda@sggw.edu.pl\\
$^{2}$ \quad Department of Clinical Engineering, Academy of Silesia in Katowice, 40-555 Katowice, Poland; bartosz.lewandowski@akademiaslaska.pl (B.L.); \hl{ewaryst.tkacz@akademiaslaska.pl (E.T.)} %MDPI: We added this email address here according to the one submitted online at susy.mdpi.com. Please confirm.
%I confirm
\\
$^{3}$ \quad \hl{Polish Honey Chamber Association,} %MDPI: Please check and confirm if the address is complete.  %I confirm
 87-100 Toruń, Poland; przemyslaw@rujna.pl\\
$^{4}$ \quad Institute of Management and Quality Sciences, Faculty of Economics, University of Warmia and Mazury in Olsztyn, 10-719 Olsztyn, Poland\\
$^{5}$ \quad Department of Commodity Science and Food Analysis, Faculty of Food Sciences, University of Warmia and Mazury in Olsztyn, 10-719 Olsztyn, Poland; renap@uwm.edu.pl\\

}


\corres{Correspondence: katarzyna.banach@uwm.edu.pl}

\abstract{\highlighting{Honey} %MDPI: Both genus and species name should be italic and protein name should not be italic, please check and confirm if their format is correct in whole paper.
 maturity is increasingly discussed in relation to product integrity, fair trade, and the classification of immature honey production as a form of adulteration. This narrative critical review examines honey maturity using evidence from peer-reviewed microbiological, physicochemical, and metabolomic studies, combined with an analysis of international and European regulatory frameworks, including Codex Alimentarius CXS 12-1981, Council Directive 2001/110/EC, Regulation (EC) No 178/2002, and Regulation (EC) No 852/2004. Particular attention is given to the interpretation of osmophilic yeast counts, water activity (\texorpdfstring{$a_w$}{aw}), moisture content, comb cell capping, fermentation, and technological dehumidification. The reviewed evidence indicates that osmophilic yeasts are natural components of honey and that their presence, expressed as colony-forming units per gram (CFU/g), should not be treated as an independent non-compliance criterion in the absence of active fermentation. Existing honey standards define compositional and quality requirements, including moisture, hydroxymethylfurfural, enzymatic activity, and absence of fermentation or effervescence, but do not establish a honey-specific CFU/g limit for yeasts. On this basis, the review formulates a functional maturity assessment framework integrating \texorpdfstring{$a_w$}{aw}, moisture, enzymatic indicators, and metabolomic biomarkers. 
{The proposed framework is presented as a conceptual model derived from the synthesis of the existing literature and requiring further multilaboratory validation prior to adoption in official control practice.} This approach may improve proportionality in honey quality assessment and reduce the risk of misclassifying microbiologically stable honeys as immature or adulterated.
}


\keyword{honey maturity; immature honey; osmophilic yeasts; water activity; fermentation risk; honey adulteration; metabolomic biomarkers; honey regulatory standards}

%%%%%%%%%%%%%%%%%%%%%%%%%%%%%%%%%%%%%%%%%%
\begin{document}
\nolinenumbers


\section{Introduction}

Honey is one of the oldest food products in human history, and at
the same time one of the most frequently adulterated commodities
on~the modern food market. According to data from the European Food
Fraud Database (FoodFraudNetwork), honey ranks third on the list of
products with the highest adulteration risk, behind only milk
and~olive oil~\cite{apimondia_fraud_v2}. In~2023, the total volume
of~honey imports into the European Union exceeded 360,000~tonnes
at~a~value close to EUR~1~billion~\cite{cbi2024}, while the EU's own
production amounted to approximately 286,000~tonnes, covering
only 60\%~of domestic demand~\cite{pollinatorhub2025}.
The EU's structural production deficit makes the market largely
dependent on supplies from third countries, including China (35\%
of~imports), Ukraine (31\%), and Argentina (12\%).
The scale of this deficit, combined with the significant price
disparity between authentic honey and syrup substitutes, creates
economic incentives for fraudulent practices~\cite{apimondia_fraud_v2}.

In the face of these challenges, Apimondia (the International
Federation of Beekeepers' Associations) plays a key role. This
organization, uniting more than 80 national organizations and
headquartered in Rome, constitutes a global forum for cooperation
aimed at promoting scientific, technical, and economic progress
in beekeeping~\cite{apimondia_web}.

The Apimondia position published in~2025 by the Working Group
on~Adulteration of Bee Products constitutes an important voice
in~the debate on product integrity~\cite{garcia2025}.
The author identifies the systemic production of immature honey,
defined as the deliberate collection of honey before the completion
of the biological ripening process in the hive, combined with its
industrial dehydration, as a separate and serious form of adulteration,
distinct from classical practices of adulteration with sugar
syrups~\cite{garcia2025}.
The significance of this position is considerable: Apimondia, as the
world's largest federation of beekeeping organizations, has real
influence on the shaping of both international standards and control
practices applied by sanitary inspectorates in EU member states.

This article does not question the existence or seriousness of
systemic immature honey production as a phenomenon violating the
principles of fair food trade. \hl{The aim of this paper is,
however, to indicate that} %MDPI: Please confirm if the bold formatting is necessary; if not, please remove it. The following highlights are the same. Please check the whole paper!
%Bold formatting is unnecessary, removed as requested!
 the Apimondia position, perhaps
inadvertently, makes an excessive generalization, equating three
qualitatively distinct phenomena:
(1)~systemic, deliberate production of immature honey for economic
gain; (2)~early honey harvesting under conditions forced by climatic
or technological factors; and (3)~the presence of osmophilic yeasts
as natural components of honey~\cite{garcia2025,munitis1976}.
The consequence of this generalization is the creation of a
disqualifying criterion based on the presence of yeasts, which
remains in contradiction with both applicable regulatory standards
and the current state of microbiological
knowledge~\cite{codex_stan12, rodriguezmachado2024}.

Analysis of applicable international standards, including Codex
Alimentarius CODEX STAN 12-1981, EU Directive 2001/110/EC, and the
draft \hl{ISO/DIS 24607:2025}  %MDPI: 1. The standard should be listed in the reference part and cited when it appears in the main text (in the form of a reference number, e.g., [2]). Please revise. 2. Please provide detailed information here, and we can help you add it to the reference part and rearrange the citation order. e.g., TB 10099-2017; Code for Design of Railway Station and Terminal. China Railway Publishing House: Beijing, China, 2017.
%Thank you for the comment. The full bibliographic details for the standard are as follows: ISO/DIS 24607; Honey — Specifications. International Organization for Standardization: Geneva, Switzerland, 2025.
standard, leads to the conclusion that
none of these documents establishes a limit on yeast cell count
(CFU/g) as a criterion for assessing the quality or authenticity of
honey~\cite{codex_stan12}.
The only legal microbiological criterion is  \hl{fermentation as an
active state,} %MDPI: Please confirm if the italics are necessary; if not, please remove them. The following highlights are the same. Please check the whole text.
%Italic is unnecessary - removed.
 not merely the presence of microorganisms capable of
fermentation~\cite{codex_stan12}. 

{
    Meanwhile, the fermentative activity of osmophilic yeasts, including \textit{Zygosaccharomyces rouxii} and \textit{Z. richteri}, is effectively inhibited at water activity $a_w < 0.60$. In honey, this value is typically, but not invariably, associated with moisture content $\leq $ $20\%$, depending on sugar composition and other matrix-specific factors~\cite{zamora2006,ruegg1981}.
}


The above premises lead to the three-part thesis of this article:
{
    first, the presence of osmophilic yeasts in honey should be interpreted in relation to water activity and the absence of active fermentation; when $a_w < 0.60$ is maintained, yeast growth and fermentation are effectively inhibited under standard storage conditions, whereas moisture content $\leq$ $20\%$ should be treated as an important regulatory criterion but not as a strict surrogate for $a_w$;
}
 second, no applicable national or international
standard penalizes the mere presence of yeasts without active
fermentation, and therefore the determination of their count by the
plate method (CFU/g) alone constitutes neither a legal nor a scientific
basis for disqualifying the product from trade; third, the concept of
\hl{functional maturity}, understood as integrated microbiological,
%Italic is unnecessary - removed.
chemical, and biological stability, constitutes a more precise and
operationalizable tool for honey quality assessment than the sole
morphological criterion of comb cell
capping~\cite{zhang2021,guo2020,sun2021}.

The article is a review and covers the following areas:
\begin{enumerate}
    \item Analysis of international regulatory frameworks concerning the microbiological parameters of honey; 
    \item Biology and metabolomics of the honey ripening process as a continuum;
    \item \textls[-15]{Microbiology of osmophilic yeasts in the context of water activity and} \mbox{fermentation risk;} 
    \item Diversity of global beekeeping production systems and their implications for standardization;     
    \item {Analysis of the EU honey market and the regulatory consequences of overly strict interpretations;}
    \item {A proposed framework for functional maturity assessment with recommendations for standardization bodies and Apimondia.}
\end{enumerate}

This article adopts a narrative critical review approach. It synthesizes peer-reviewed scientific literature and selected regulatory, standardization, and institutional documents relevant to honey maturity. 
{Peer-reviewed literature was identified through searches of PubMed, Web of Science, and Scopus, using combinations of the following search terms: \hl{honey maturity}, \hl{immature honey}, \hl{osmophilic yeasts}, \hl{water activity honey}, \hl{honey fermentation}, \hl{honey adulteration}, \hl{honey metabolomics}, \hl{honey ripening}, and \hl{Zygosaccharomyces honey}. The database search covered publications from 1980 to March 2026 and was not restricted by language at the query stage; however, the final synthesis included only English-language sources, as these could be reliably assessed by all authors. Regulatory and institutional documents were selected directly from official sources: FAO/WHO (Codex Alimentarius), the European Commission, ISO technical committees, national standards bodies, and major international beekeeping federations. A source was included if it directly addressed honey quality criteria, microbiological standards, or maturity assessment methodology.} Particular attention is given to osmophilic yeasts, water activity, moisture content, fermentation risk, technological dehumidification, and honey adulteration. The reviewed sources include studies on honey microbiology, physicochemical and enzymatic indicators of ripening, metabolomic markers, international honey standards, and EU food law. The analysis is organized around four thematic areas: the regulatory interpretation of honey quality and fermentation, the biological and physicochemical mechanisms of honey ripening, the relationship between yeast counts and fermentation risk, and the applicability of maturity criteria under different climatic and production systems.

% ============================================================

\section{International Regulatory Frameworks: What Standards \linebreak  Actually Prohibit}\label{2}

The assessment of honey quality and authenticity is based on a network
of interrelated normative documents of both global and regional scope.
Understanding their scope and limitations is crucial for evaluating
whether the presence of osmophilic yeasts can constitute an independent
basis for disqualifying a product from commercial trade.

\subsection{Codex Alimentarius, CODEX STAN 12-1981}

The primary document regulating honey quality at the international
level is the Codex Alimentarius Standard for Honey
({CODEX STAN 12-1981,}  last revision 2023), developed jointly
by FAO and WHO~\cite{codex_stan12}. This standard defines honey
as a natural product made by honey bees from the nectar of plants
or from secretions of living parts of plants or excretions of
plant-sucking insects on the living parts of plants, which the bees
collect, transform by combining with specific substances of their own,
deposit, dehydrate, store, and leave in the honeycomb to ripen and
mature~\cite{codex_stan12}.

In terms of physicochemical parameters, CODEX STAN 12-1981 establishes
the following quantitative limits: water content $\leq$ $20.0\%$
(for heather honey \textit{{Calluna}} $\leq$ $23.0\%$), content
of~hydroxymethylfurfural (HMF) $\leq$ $40$~mg/kg (or $\leq$$80$~mg/kg
for honeys from countries with a tropical climate), diastase
activity $\geq$ $8$~Schade units (or $\geq$$3$~Schade units for honeys
with naturally low enzymatic activity, such as acacia honey
\textit{Robinia pseudoacacia}), and reducing sugar
content~\cite{codex_stan12}. The standard contains \hl{no 
 limit
on yeast count expressed in CFU/g}. The only microbiological criterion
is the prohibition on trading honey \hl{``that has begun to ferment
or effervesce''}~\cite{codex_stan12}. This is a distinction of
fundamental legal importance: the disqualifying criterion is
\hl{the state of active fermentation}, not \hl{the potential
cause} (the presence of microorganisms capable of fermentation).

\subsection{EU Directive 2001/110/EC and National Law}

At the European Union level, the primary legal act regulating honey
trade is Council Directive 2001/110/EC of 20 December 2001 on honey,
implemented in the Polish legal order by a regulation of the Minister
of Agriculture and Rural Development~\cite{eudir2001110}.
This Directive, modeled on the Codex standard, does not establish
a~CFU/g limit for yeasts or fungi in honey. Microbiological parameters
of honey are subject to general food hygiene provisions, in particular
\hl{Regulation (EC) No.~852/2004} %MDPI: 1. The standard should be listed in the reference part and cited when it appears in the main text (in the form of a reference number, e.g., [2]). Please revise. 2. Please provide detailed information here, and we can help you add it to the reference part and rearrange the citation order. e.g., TB 10099-2017; Code for Design of Railway Station and Terminal. China Railway Publishing House: Beijing, China, 2017.
%Thank you for the comment. The full bibliographic details are as follows: Regulation (EC) No 852/2004 of the European Parliament and of the Council of 29 April 2004 on the hygiene of foodstuffs. Official Journal of the European Union, L 139, 30 4.2004, pp. 1–54.
  of the European Parliament and of the
Council on the hygiene of foodstuffs~\cite{eudir2001110}.

In the control practice of veterinary and sanitary inspection
laboratories, there is an unofficial hygiene threshold of
$100$~CFU/g for the total count of fungi and yeasts, used as a general
indicator of production process hygiene; however, this threshold is
not grounded in specific provisions regarding honey and cannot
constitute an independent basis for declaring the product
non-compliant with the requirements of food law within the meaning
of Article~14 of Regulation~(EC)~No.~178/2002~\cite{eudir2001110}.
Directive 2001/110/EC also permits the possibility of filtering,
heating (to specified temperatures), and blending batches of honey
as legal technological treatments, provided they do not lead to
a~change in its fundamental character~\cite{eudir2001110}.


\subsection{United States Pharmacopeia (USP) and Other National Standards}

The United States Pharmacopeia (USP) classifies honey as an excipient
for pharmaceutical and food use. USP requirements include
physicochemical parameters similar to those of Codex but also do not
contain a CFU/g limit for osmophilic yeasts~\cite{tsagkaris2021}.
A similar approach is taken by national standards applicable in the
major honey-exporting countries---China (\hl{GB/T~18796-2012}%MDPI: 1. The standard should be listed in the reference part and cited when it appears in the main text (in the form of a reference number, e.g., [2]). Please revise. 2. Please provide detailed information here, and we can help you add it to the reference part and rearrange the citation order. e.g., TB 10099-2017; Code for Design of Railway Station and Terminal. China Railway Publishing House: Beijing, China, 2017.
% Thank you for the comment. We apologize for the typo — the correct standard number is GH/T 18796-2012 (not GB/T 18796-2012). The full bibliographic details are as follows: GH/T 18796-2012; Honey. Standards Press of China: Beijing, China, 2012.
), Argentina
(Código Alimentario Argentino, CAA), and Brazil \hl{(IN MAPA No.~89/2000)}%MDPI: 1. The standard should be listed in the reference part and cited when it appears in the main text (in the form of a reference number, e.g., [2]). Please revise. 2. Please provide detailed information here, and we can help you add it to the reference part and rearrange the citation order. e.g., TB 10099-2017; Code for Design of Railway Station and Terminal. China Railway Publishing House: Beijing, China, 2017.
%Thank you for the comment. We apologize for the error — the correct reference is Instrução Normativa No. 11/2000 (not No. 89/2000). The full bibliographic details are as follows: Brazil, Ministério da Agricultura e Abastecimento. Instrução Normativa No. 11, de 20 de outubro de 2000. Regulamento Técnico de Identidade e Qualidade do Mel. Diário Oficial da União: Brasília, Brazil, 2000.
---which define honey quality by chemical parameters and enzymatic
activity, not by yeast count~\cite{tsagkaris2021}.

\subsection{Comparative Overview of Standards}

Table~\ref{tab:standardy} presents a comparison of key honey quality
parameters in applicable international and regional standards with
respect to microbiological criteria.

\begin{table}[H]




\caption{Comparison of honey quality parameters in major international
         standards. Abbreviations: n.a., not applicable (no limit
         in the document); CFU, colony-forming units;
         HMF, \mbox{hydroxymethylfurfural.}}
\label{tab:standardy}


\begin{adjustwidth}{-\extralength}{0cm}


\begin{tabularx}{\fulllength}{cCCCcc}
\toprule
\textbf{Standard} &
\textbf{Moisture} &
\textbf{HMF} &
\textbf{Diastase} &
\textbf{CFU/g Limit} &
\textbf{Criterion} \\
 & \textbf{(\%)} & \textbf{(mg/kg)} & \textbf{(Schade u.)} & \textbf{Yeasts} & \textbf{Microbiolog.} \\
\midrule
{Codex STAN 12-1981}        & $\leq$$20.0$ & $\leq$$40$  & $\geq$$8$   & \hl{none} & prohibition of fermentation \\%MDPI: \hl{} %MDPI: Please add an explanation for the use of bold in the table footer. If the bold is unnecessary, please remove it. The following highlights are the same.
%Bold removed

{EU Directive 2001/110/EC} & $\leq$$20.0$ & $\leq$$40$  & $\geq$$8$   & \hl{none} & prohibition of fermentation \\
{USP (honey monograph)}    & $\leq$$21.0$ & n.a.       & n.a.        & \hl{none} & prohibition of fermentation \\
{Reg. (EC) 852/2004}       & n.a.          & n.a.       & n.a.        & \hl{none}  & general food $^{*}$ \\
\bottomrule

\end{tabularx}

\begin{adjustwidth}{--\extralength}{0cm}
\noindent\footnotesize $^{*}$~The 100 CFU/g value may be used in some laboratory or inspection practices as a general hygiene indicator for yeasts and molds; however, it is not a honey-specific legal criterion established by Regulation (EC) No 852/2004.
\end{adjustwidth}
\end{adjustwidth}

\end{table}



\subsection{Regulatory Conclusion}

Analysis of the above documents leads to an unequivocal conclusion:
no applicable national or international standard treats \hl{the
presence of osmophilic yeasts per~se} as a basis for disqualifying
honey from commercial trade. The disqualifying criterion is exclusively
the \hl{biological outcome} (active fermentation and effervescence),
not the \hl{potential cause} (the presence of microorganisms capable
of fermentation). 

{This distinction is of fundamental practical importance: honey in which osmophilic yeasts are present at any measured count should not be disqualified on this basis alone, provided that water activity $a_w < 0.60$ is maintained, no signs of active fermentation are observed, and other compositional criteria, including the applicable moisture limit, are met}~\cite{codex_stan12,eudir2001110}.

A CFU/g measurement alone, without a simultaneous assessment of water
activity ($a_w$), temperature, and storage time, is methodologically
insufficient for declaring the product non-compliant with the
requirements of food law~\cite{tsagkaris2021,schoder2026}.






\section{Biology of Honey Ripening: A Continuous Process, Not a Binary Event}

Honey ripening is a complex biotransformation process in which
nectar or honeydew undergoes a series of enzymatic, physicochemical,
and microbiological transformations before reaching a state of
qualitative stability. The widespread equating of honey maturity
with the fact of comb cell capping with wax---common in beekeeping
practice and veterinary control---constitutes a simplification that
is inadequate from a biochemical perspective. Capping is one of the
final, morphological indicators of ripening, not synonymous with the
achievement of full chemical and microbiological
stability~\cite{zhang2021,sun2021}.

\subsection{Enzymatic Transformation of Nectar}
The honey production process is initiated at the very moment nectar
or honeydew is collected by forager bees. 

Nectar, dominated by sucrose and containing a high water content, undergoes the first enzymatic transformation in the honey crop of the insect, where it mixes with secretions from the bee’s hypopharyngeal glands~\cite{biochemreactions2022}. The key enzymes initiating
the \mbox{transformation are:}

\begin{itemize}
    \item \hl{Invertase} %MDPI: Please confirm if the bold formatting is necessary; if not, please remove it. The following highlights are the same. %Bold removed!
 ($\beta$-D-fructosidase): It catalyzes
          the hydrolysis of sucrose into fructose and glucose;
          produced by the hypopharyngeal glands of worker
          bees~\cite{biochemreactions2022};
          its activity increases gradually from low values
          in immature honey to maximum values in mature honey,
          making it one of the most sensitive indicators of
          the degree of maturity~\cite{zhang2021}.
    \item \hl{Glucose oxidase}: It oxidizes glucose to gluconic
          acid with the release of hydrogen peroxide (H$_2$O$_2$);
          gluconic acid is responsible for lowering honey pH to the
          range of 3.5--4.5 and constitutes the dominant organic
          acid in mature honey; H$_2$O$_2$ serves as a natural
          antibacterial agent~\cite{biochemreactions2022,zhang2021}.
    \item \hl{Diastase} 
 ($\alpha$- and $\beta$-amylase):
          It catalyzes the hydrolysis of starch and nectar
          polysaccharides into glucose and maltose; diastase
          activity is a parameter standardized by Codex Alimentarius
          ($\geq$$8$~Schade units) and serves as an indirect indicator
          of thermal integrity and honey
          maturity~\cite{codex_stan12,biochemreactions2022,zhang2021}.
    \item \hl{Catalase}: 
 It decomposes excess H$_2$O$_2$,
          regulating the balance between antibacterial activity
          and protection of the remaining honey components
          from oxidation~\cite{biochemreactions2022}.
\end{itemize}

Studies of enzymatic activity at five successive stages of
transformation---from floral nectar, through the honey crop of
foragers and house bees, to uncapped and capped hive cells---showed
a gradual increase in invertase, amylase, and glucose oxidase
activity at each successive stage of
processing~\cite{enzymesrole2023}. These data confirm that
enzymatic transformation proceeds as a \hl{continuum} along
the entire production chain, not as a step-change event occurring
at the moment of cell capping~\cite{enzymesrole2023,zhang2021}.

In parallel with enzymatic transformations, house bees actively
dehydrate honey through hive ventilation: wing beats create a
constant airflow over open comb cells, accelerating water
evaporation and reducing the moisture of the raw material from
values characteristic of nectar to a final approximately 18.5\%
in honey before capping~\cite{zhang2021}. Full maturation under
natural conditions typically takes several days to a few weeks,
with this time potentially being significantly extended under
conditions of high relative
\mbox{humidity~\cite{zhang2021,tadele2025}.}


\subsection{Ripening as a Process: Metabolomic Evidence}

The thesis of the binary nature of honey maturity---according to
which honey is either immature (uncapped) or mature (capped)---finds increasingly weak support in current metabolomic research.
Sun et al.~\cite{sun2021} in a study of rapeseed honey
(\textit{Brassica napus}) demonstrated that maturity indicators,
including water content, fructose-to-glucose ratio, invertase and
diastase activity, and organic acid content, change {\hl{gradually}
during the ripening process and reach stable values independently
of the moment of cell capping \mbox{by bees.}

Guo et al.~\cite{guo2020} conducted a comprehensive comparative
analysis of the chemical profile and biological activity of mature
and immature honeys using HPLC/QTOF/MS-based metabolomics. The
authors identified significant differences in polyphenol content,
antioxidant activity (DPPH, FRAP), and amino acid profile between
honey collected in an immature state and mature honey; however,
these differences were closely correlated with chemical parameters
(water content and enzymatic activity) rather than solely with
the fact of capping~\cite{guo2020}. Importantly, some samples of
honey collected before full capping exhibited a metabolomic profile
similar to mature honey, provided they met the criterion of low
moisture, which constitutes direct confirmation of the concept
of functional maturity.

Sun et al.~\cite{sun2021} using GC-MS and LC-MS identified
decenedioic acid, a fatty acid of bee origin, as a potential
molecular biomarker differentiating mature from immature honey.
This biomarker, validated for rapeseed, acacia, and jujube honeys,
was significantly more abundant in mature honey ($p < 0.01$),
confirming that ripening is a process of gradual accumulation of
specific metabolites, not a binary transition~\cite{sun2021}.


{Beyond fatty acid biomarkers, the metabolomic fingerprint of honey maturity encompasses several additional compound classes. Polyamines, particularly spermidine and spermine of plant and bee origin, accumulate progressively during ripening and have been proposed as markers of enzymatic activity and cellular integrity in the maturing product~\cite{guo2020,biochemreactions2022}. Flavonoids (quercetin, kaempferol, and luteolin) and phenolic glycosides derived from floral nectar undergo concentration and partial transformation during the ripening process; their profiles are correlated with diastase and invertase activity and constitute part of the antioxidant capacity indexed by DPPH and FRAP assays~\cite{guo2020}. Organic acids, including gluconic, citric, malic, and succinic acids, increase in concentration as glucose oxidase activity rises during ripening, providing both a pH indicator and a chemical maturity marker independent of moisture content~\cite{zhang2021,biochemreactions2022}. These compound classes have been identified in multiple botanical types and are analytically accessible by LC-MS and NMR-based metabolomics platforms~\cite{nmrhoney2023}, supporting their potential inclusion in a multiparametric maturity assessment framework.}

It is worth emphasizing that exposure of honey to high temperatures,
even after full capping, can cause enzyme degradation and effectively
revert quality parameters to values characteristic of immature
honey~\cite{biochemreactions2022}.
Capped honey subjected to inappropriate heat treatment may exhibit
diastase activity below the Codex threshold (<8~Schade units)
and elevated HMF content, demonstrating that the morphological
criterion of capping does not guarantee the maintenance of
functional maturity throughout the supply chain.

This is also confirmed by beekeeping practice, which treats the
fact of cell capping merely as an indicator of potential honey
maturity. On one hand, it sometimes happens that under certain
conditions bees cap comb cells in which honey is not yet mature
(water content exceeds 20\% and shortly after extraction the honey
ferments). On the other hand, common practice involves harvesting
honey from partially capped frames (usually with a threshold of
three-quarters of cells capped in the frame), provided that in
the uncapped cells the honey is mature. The actual indicator of
honey maturity is the water content, which the beekeeper determines
either by refractometry or by other tests (e.g., shaking the frame),
whose result is directly dependent on water content.

\subsection{The Concept of Functional Maturity}\label{3.3}
Based on the above review, functional honey maturity can be operationally defined as an integrated product state characterized by the simultaneous fulfillment of three groups of criteria:

\begin{enumerate}
    \item \hl{Microbiological stability:} %MDPI: Please confirm if the bold formatting is necessary; if not, please remove it. The following highlights are the same. %Bold removed.
    absence of active
          fermentation; water activity $a_w < 0.60$; moisture
          $\leq$ $20\%$; osmophilic yeast count below the growth
          threshold under given conditions of temperature
          and storage~\cite{zamora2006}.

    \item \hl{Chemical stability:} fructose-to-glucose ratio
          (F/G) $\geq$ $0.9$; HMF content $\leq$ $40$~mg/kg;
          diastase activity $\geq$ $8$~Schade units; invertase
          activity $\geq$ $4$~units; gluconic acid content and
          organic acid profile within the reference range
          for the given botanical type~\cite{codex_stan12,zhang2021}.

    \item \hl{Biological stability:} antioxidant activity
          (DPPH, FRAP) at the level appropriate for the given
          botanical and geographical type; polyphenol and flavonoid
          content not reduced relative to reference values;
          antibacterial activity (assessed by H$_2$O$_2$ production
          and the presence of defensins)
          unchanged~\cite{guo2020}.
\end{enumerate}

This definition deliberately separates the morphological criterion
(cell capping) from the quality criterion, recognizing that honey
may achieve functional maturity before full comb capping (e.g., at
18\% moisture in a temperate climate). This concept is consistent
with the Codex regulatory approach, which, as demonstrated in
Section~\ref{2}, penalizes the state of fermentation, not
morphology~\cite{codex_stan12,sun2021}.






\section{Microbiology of Yeasts in Honey:
Presence vs.\ Activity vs.\ Adulteration}\label{4}

The microbiology of honey is a field relatively rarely analyzed
in the context of product authenticity assessment, despite the fact
that microorganisms---among them osmophilic yeasts---are an
inseparable component of every honey produced under natural
conditions. This section states and justifies the central
microbiological thesis of the article: the mere presence of
osmophilic yeasts in honey is a biologically normal phenomenon,
and does not constitute grounds for declaring the product immature,
let alone adulterated, provided that water activity ($a_w$) and
water content are maintained at levels prescribed by applicable
standards.


\subsection{Osmophilic Yeasts as a Natural Component of Honey}


\subsubsection{Biological Characteristics of Osmophilic Yeasts}

Osmophilic yeasts, also called osmotolerant or xerotolerant, are
microorganisms capable of growth and metabolism in environments
with very high osmotic pressure, i.e., at water activity $a_w$
close to $0.60$--$0.62$~\cite{ruegg1981,waaw2006}. In honey,
the dominant role is played primarily by species of the genus
\textit{Zygosaccharomyces}:
\textit{Zygosaccharomyces rouxii} (Boutroux) Yarrow, first described
in 1893 as yeasts isolated from sucrose solutions, constitutes the
model species for fungal osmotolerance. Its minimum water activity
for growth is $a_w = 0.62$, making it one of the most resistant
osmotolerant fungi known to
science~\cite{waaw2006}. Under conditions below this threshold,
\textit{Z.~rouxii} remains in a state of cryptobiosis; cells
retain viability and germination capacity but show no metabolic
activity or growth~\cite{waaw2006,ruegg1981}.

More recent studies using culture-dependent and culture-independent
methods (metagenomics) have expanded the list of yeasts isolated
from honey to include further unconventional species, including
\textit{Candida} spp., \textit{Rhodosporidiobolus ruineniae},
\textit{Clavispora lusitaniae}, and \textit{Metschnikowia
chrysoperlae}~\cite{rodriguezmachado2024}.
Metagenomic studies indicate that \textit{Zygosaccharomyces}
remains the dominant genus in both western honey bee
(\textit{Apis mellifera}) honeys and eastern honey bee
(\textit{Apis cerana}) honeys, regardless of country of
origin~\cite{yeastbiodiversity2025}.

\subsubsection{Sources of Colonization and Prevalence}
Osmophilic yeasts are an integral component of the hive ecosystem
and beekeeping environment, which not only precludes the
justifiability of imposing restrictions in this regard but also
excludes the possibility of their elimination from honey without
applying extreme technological interventions. The main colonization
routes include floral nectar, in which yeast counts are typically
$10^2$--$10^4$~CFU/ml depending on the plant species, microclimate,
and season; pollen and bee bread, which are carriers of active
yeast and lactic acid bacteria populations; as well as the bee's
digestive tract---the honey crop and intestine constitute a natural
reservoir of microorganisms~\cite{alwaili2012}.
Additional sources include soil, wax, and hive surfaces, whose yeast
biofilm is difficult to eliminate without disinfectants incompatible
with organic production~\cite{alwaili2012}.

The presence of osmophilic yeasts in honey is therefore an inevitable
phenomenon resulting from the biology of the production process,
not from production deficiencies. This is confirmed by the literature
from various climatic and geographical zones, encompassing both
western and eastern honey bee
honeys~\cite{yeastbiodiversity2025}.
The volume of yeasts found in normal non-fermenting commercial
honey typically ranges from several dozen to several hundred CFU/g;
values above 1000~CFU/g are recorded in correct samples without
signs of fermentation, provided moisture does not exceed
$18\%$~\cite{analytica2020}.
{To provide concrete empirical context, Ruegg and Blanc~\cite{ruegg1981}
reported osmophilic yeast counts of $10$--$10^3$~CFU/g in commercially
traded European honeys showing no fermentation signs, and Rodrigue\-z-Machado
et al.~\cite{rodriguezmachado2024} identified counts ranging from below
detection limits to $1.5 \times 10^3$~CFU/g across 120 honey samples
of diverse botanical origin, none of which showed active fermentation
at moisture $\leq$ $18\%$. Metagenomic surveys confirm that
\textit{Zygosaccharomyces} populations in this CFU/g range represent
a stable, metabolically inactive reservoir under $a_w < 0.60$
conditions~\cite{yeastbiodiversity2025}.}



\subsection{Water Activity as a Critical Control Parameter}


\subsubsection{Definition and Significance of Water Activity in Honey}

Water activity ($a_w$) is a thermodynamic measure of the
availability of water in a food product for chemical reactions
and microbiological activity, defined as the ratio of the vapor
pressure above the product to the vapor pressure of pure water
at saturation at the same
temperature~\cite{waaw2006}.
In contrast to water content (moisture), expressed as mass
percentage, $a_w$ reflects the actual availability of water to
microorganisms and is the primary parameter determining microbial
growth, enzymatic activity, and non-enzymatic browning of the
product~\cite{waaw2006}.

Honey is a supersaturated solution of simple sugars, fructose
($\sim$$38\%$) and glucose ($\sim$$30\%$), and higher
oligosaccharides, in an environment of minimal water content.
The high osmolarity of the solution resulting from this saturation
causes water molecules to be strongly bound to sugar molecules,
and their thermodynamic activity is drastically
reduced~\cite{zamora2006,waaw2006}. In honey with moisture
$\leq$ $20\%$, water activity is typically $a_w < 0.60$,
although the exact value depends on sugar composition, water
content, and measurement temperature~\cite{ruegg1981,gleiter2006}.
The study of Ruegg and Blanc~\cite{ruegg1981}, cited by all
subsequent food microbiology textbooks as a fundamental reference
for honey microbiology, demonstrated that for the moisture range
of 16--21\% water activity falls between $a_w = 0.56$ and
$a_w = 0.63$, with the value $a_w = 0.60$ corresponding to
moisture of $\approx$$17.5$--$18.5\%$ depending on sugar
composition.

\subsubsection{Osmophilic Yeast Growth Threshold and Water Activity}

The minimum water activity required for osmophilic yeast growth
is $a_w = 0.61$--$0.62$. 
      Below this threshold, yeast growth is effectively inhibited under standard storage conditions, and cells may remain viable but metabolically inactive or markedly restricted in activity~\cite{ruegg1981,alwaili2012}. In the context of honey, this means that in products with $a_w < 0.60$, osmophilic yeasts are not expected to grow or initiate fermentation under normal commercial storage conditions, irrespective of the measured CFU/g count.


The practical implication is direct: 
honey containing even $10^4$~CFU/g of osmophilic yeasts but with moisture of approximately 17\% and ($a_w \approx 0.57$--$0.58$) should remain microbiologically stable under standard storage conditions~\cite{alwaili2012}.


The Analytica Laboratories document formulates this principle
as a multiparametric model: fermentation is highly probable only
at moisture > 19$\%$ \hl{and} yeast count > 10~CFU/g
simultaneously~\cite{analytica2020}, which \hl{a~contrario} %MDPI: Please confirm if the italics are necessary; if not, please remove them. The following highlights are the same.
precludes product disqualification on the basis of the CFU/g
parameter alone.

For reasons of microbiological stability, water activity is an
incomparably more important parameter than the percentage water
content. It should be noted that for certain types of honey
(e.g., heather honey), legal requirements permit water content
levels above 20\%. This is due to the fact that the specific
composition of these honeys ensures microbiological stability at
water content exceeding 20\%, and these honeys should simultaneously
be considered ``mature'' at the values specified in legal
requirements. If water activity studies demonstrate that the water
activity of these honeys does not exceed the threshold
($a_w < 0.60$), this parameter would constitute a significantly
more universal parameter for determining the microbiological
stability and degree of maturity of honey.

\subsubsection{Glucose Crystallization as a Dynamic Factor Modifying Water Activity}

The natural process of glucose crystallization in honey is a factor
deserving particular attention in the context of fermentation risk,
as it constitutes a mechanism by which initially microbiologically
stable honey can become susceptible to fermentation without any
external intervention. During crystallization, glucose transitions
from the liquid phase to the solid phase as glucose monohydrate
(C$_6$H$_{12}$O$_6 \cdot$ H$_2$O). Each glucose molecule binds
one water molecule, effectively removing it from the liquid
phase~\cite{zamora2006}. However, the liquid phase remaining
after glucose crystallization becomes enriched in fructose and
water, causing an \hl{increase} in water activity in the
remaining liquid phase, even though the total water content in
the sample does not change~\cite{zamora2006,ijfans2019}.

Zamora and Chirife~\cite{zamora2006} in a study of 50 samples
of crystallized and re-dissolved Argentine honeys demonstrated
that crystallization causes an increase in the $a_w$ of the liquid
phase by $0.03$--$0.04$ units compared to honey in the liquid
state. For honeys with initial moisture of $17.5$--$18.0\%$
($a_w \approx 0.58$--$0.59$), crystallization can shift the
$a_w$ of the liquid phase above the threshold of $0.62$, thereby
activating the osmophilic yeast population present in the
product~\cite{zamora2006,ijfans2019}.
This phenomenon leads to a key methodological conclusion: the
microbiological control of honey must take into account the phase
state of the product (liquid vs.\ crystallized) since the same
CFU/g count carries different fermentation risk depending on whether
the honey is liquid or crystallized. These processes further support the need to interpret yeast counts together with water activity and the physical state \mbox{of honey.}

Given that honey crystallization proceeds in a differentiated
manner (i.e., through different crystal forms with regard to both
structure and dynamics of formation), changes in water activity
and thus the potential for yeast proliferation in crystallizing
honey are variable. This variability is observed not only with
respect to the time/progress of crystallization but also the
sampling point within the container. Considering that crystallization
propagation occurs at the molecular level, the greatest potential
for water concentration occurs in the pericrystalline and
intercrystalline spaces. Mainly through diffusion, this fluctuation
tends to equalize in areas increasingly distant from the crystal
surface. It should be noted that the interphase equilibrium is a
dynamic process in which the direction of phase transitions is
determined by temperature (variable storage conditions). Depending
on the conditions accompanying these phase transitions (such as
temperature, solution density, oxygen availability, honey aeration,
the free spaces formed in honey during crystallization, the rate of
crystallization resulting from the proportions of individual sugars,
etc.), the potential for yeast proliferation in solution is dynamic
and varies throughout the volume of honey. A CFU/g measurement
without precise specification of the phase state and~$a_w$ is
therefore methodologically incomplete.

From the above it follows that if water activity could determine
honey maturity, if this parameter were specified in legal
requirements, the water activity value should be set at a level
achievable at the time of honey extraction, and should also be
maintained (without further technical measures) despite differences
in temperature and phase transitions occurring during storage and
distribution.


\subsection{Yeast Count as a Risk Indicator, Not a Disqualifying Parameter}\label{4.3}


\subsubsection{Multiparametric Model of Fermentation Risk}

Analysis of the scientific literature and practical documents
indicates the widespread acceptance of a multiparametric model
of honey fermentation risk, in which no single indicator is
sufficient to determine the microbiological stability or
instability of the product. Fermentation risk is a function
of at least four \hl{variables:} %MDPI: 1. The format of the variables in the formulas should be consistent both in text and in equations. Please check, confirm and modify it in the whole text. 2. Please check all formulas carefully to make sure all formulas are correct and there are no duplicate formulas.

%1. Checked throughout the text. One inconsistency found and corrected:
%   line 335, $a_w < 0.60\%$ changed to $a_w < 0.60$ (water activity  is dimensionless; the % sign was erroneous). All other variable formats are consistent between text and equations.
%2. The document contains one numbered equation (eq:fermentation_risk).
%   No duplicate formulas found.

\begin{equation}
R_f = f\!\left(a_w,\; T,\; t,\; N_{CFU}\right)
\label{eq:fermentation_risk}
\end{equation}
 where $R_f$ is fermentation risk (probability of initiating
active fermentation), $a_w$ is water activity, $T$ is storage
temperature [$^\circ$C], $t$ is storage time [days],
and $N_{CFU}$ is the osmophilic yeast count [CFU/g].

According to this model, yeasts can ferment honey only when
{all} growth-promoting parameters are simultaneously fulfilled:
$a_w$ must exceed the threshold of $0.62$, temperature must fall
within the optimal range for yeasts ($10$--$27\,^\circ$C),
and the exposure time must be sufficient to reach the cell density
initiating noticeable fermentation~\cite{apinz2021}.
The New Zealand Beekeeping Industry document states directly that
honey with moisture > 19\% and yeast count > 10~CFU/g is
susceptible to fermentation, which \hl{a contrario} implies
that {
honey with moisture < 17\%, even with a substantially higher yeast count, is not expected to ferment under normal commercial storage conditions}~\cite{apinz2021}.

Table~\ref{tab:fermentation_risk} illustrates the matrix model
of fermentation risk as a function of moisture and yeast count,
based on available empirical data.

\begin{table}[H]

\caption{Estimated honey fermentation risk as a function
         of moisture and osmophilic yeast count
         (storage temperature $10$--$27\,^\circ$C).
         Source: own compilation based on~\cite{analytica2020,apinz2021}.}
\label{tab:fermentation_risk}


\begin{adjustwidth}{-\extralength}{0cm}
\begin{tabularx}{\fulllength}{l *{3}{>{\centering\arraybackslash}X}}
\toprule
\multirow{2.5}{*}{\textbf{Moisture [\%]}} &
\multicolumn{3}{c}{\textbf{Yeast Count [CFU/g]}} \\
\cmidrule{2-4}
 & \textbf{<10} & \textbf{10--1000} & \textbf{>1000} \\
\midrule
$\leq$17.0       & {negligible risk}       & {negligible risk}               & {negligible risk} \\
$17.1$--$18.0$  & {negligible risk}       & low risk              & moderate risk \\
$18.1$--$19.0$  & {negligible risk}       & moderate risk         & high risk \\
$19.1$--$20.0$  & low risk      & high risk             & {very high risk} \\
$>$20.0          & moderate risk & {high risk} & {very high risk} \\
\bottomrule
\end{tabularx}
\end{adjustwidth}

\end{table}




{The risk categories presented in Table~\ref{tab:fermentation_risk} are heuristic estimates compiled from two practitioner-oriented documents~\cite{analytica2020,apinz2021} and represent the consensus of applied honey quality practice. They have not been derived from a controlled experimental study under defined storage conditions; the term ``no risk'' should be understood as indicating negligible probability of fermentation initiation under normal commercial storage, not an absolute thermodynamic guarantee. The table is intended as a decision-support tool and should be interpreted alongside direct $a_w$ measurement rather than used as a standalone \mbox{classification system.}}

\subsubsection{Limitations of the Plate Method (CFU/g) as a Risk Assessment Tool}


The standard enumeration of yeasts and molds in foods is commonly based on plating serial dilutions on selective media, followed by incubation and colony counting.


This method, while widely used and well-validated, has a number of significant
limitations in the context of honey fermentation risk assessment:


\begin{enumerate}
    \item \hl{Failure to distinguish viable from dead cells:}
          the plate method counts only cells capable of forming
          colonies on solid medium under incubation conditions;
          living cells in a state of deep anabiosis induced by
          low honey $a_w$ may fail to grow on the medium or grow
          atypically, leading to underestimation of the true yeast
          count;
    \item \highlighting{Lack of correlation with in situ \boldmath{$a_w$:}} the
          CFU/g result is a measure of the potential population size
          but does not indicate whether under the conditions prevailing
          in the tested product this population is active;
          only the combination of CFU/g with $a_w$ measurement allows
          assessment of the actual fermentation
          risk~\cite{ruegg1981};
    \item \hl{Lack of species specificity:} standard selective
          media do not differentiate yeast species with different
          fermentation potential; \textit{Metschnikowia}
          spp.\ with low fermentation potential are counted
          identically to \textit{Z.~rouxii} with high
          potential~\cite{rodriguezmachado2024}.
\end{enumerate}

\textbf{\hl{Legal--scientific conclusion:} %MDPI: Please check if this paragraph should be a part of above list content.
} in light of the above
methodological limitations, the mere determination of yeast count
(CFU/g) in honey, without a simultaneous measurement of $a_w$,
determination of the product's phase state, storage temperature,
and time since harvest, does not constitute a sufficient scientific
basis for assessing fermentation risk, let alone for disqualifying
the product from commercial trade. Given that no applicable standard
(Codex, EU, ISO) establishes a CFU/g limit for honey
(cf.~Section~\ref{2}), a potential inspector's decision to disqualify
a product based solely on CFU/g count would be devoid of both
scientific and legal grounds.


\subsection{Fermentation as a Dynamic and Post-Harvest Event}


The Apimondia position and much of the literature on honey
adulteration tacitly assume that fermentation detected in honey
is evidence of its immaturity at the time of harvest.
This assumption is, however, incorrect from the perspective of
microbiological dynamics since fermentation is an event that can
occur at any stage of the supply chain, not only at the time
the product leaves the hive~\cite{apinz2021,analytica2020}.

Correctly harvested honey that meets all normative criteria,
with moisture of, for example, $18.5\%$ ($a_w \approx 0.60$)
and a yeast count of $200$~CFU/g, stored for six months in a
warehouse at a temperature of $22\,^\circ$C and relative humidity
above 70\%, may locally exhibit higher water activity
(crystallization, temperature changes, air access, etc.), shifting
the $a_w$ above the yeast growth threshold and initiating
fermentation~\cite{apinz2021}.

Border inspection, sampling imported honey after several months
of transport and storage, takes measurements at a point removed
in time and space from the moment of harvest; the determination of
fermentation or elevated active yeast count at that moment does not
allow, without analysis of temperature and humidity documentation
in the supply chain (\hl{cold chain monitoring}), reliable
reconstruction of the product's state at the time of honey
extraction (verifying the degree of capping) or departure from
the apiary~\cite{tsagkaris2021,schoder2026}.
Therefore, fermentation observed at the point of inspection should not be treated, by itself, as proof of deliberate adulteration at the production stage. It may indicate product instability or non-compliance, but attribution to intentional immature honey production requires additional evidence from the extraction, processing, storage, or \mbox{supply-chain stages.}

Deliberate adulteration consisting of
harvesting immature honey followed by technological treatments
aimed at its microbiological stabilization can only be established
through controls at the honey extraction stage, and effective
limitation may be achieved by specifying prohibited technologies
(e.g., vacuum dehydration). The non-use of prohibited technologies
should be a mandatory stage of official controls of apiaries and
establishments involved in honey extraction and placing on the
market. Neither yeast count nor fermentation permits inference
about this.


\section{Global Production Systems and Climatic Conditions}\label{5}

Honey quality standards, including the Apimondia position, were
developed on the basis of the European temperate climate and
frame beekeeping, yet more than 70\% of global honey production
comes from regions that are climatically, technologically, and
culturally different from
Europe~\cite{garcia2025,eudir2001110}. Applying uniform morphological and microbiological criteria to honeys from different production systems is methodologically limited and may create disproportionate assessment outcomes for producers operating under different climatic and technological conditions~\cite{phiri2022,tadele2025}.


\subsection{Global Landscape of Apicultural Production}


According to FAO data from 2024, the global number of managed
colonies exceeded 101.7~million, with Asia dominating
at a 44.4\% share, and Europe holding 24.7\%~\cite{destatis2026}.
Honey production in 2022 reached 1.83~million tonnes, with
Asia accounting for 48.2\%, Europe for 22.8\%, the Americas
for 18.5\%, and Africa for 8.5\%~\cite{phiri2022}. Adopting
the capping criterion as the sole determinant of maturity in
practice favors European honeys and discriminates against honeys
from tropical and traditional
systems~\cite{tadele2025,phiri2022}.





Relative humidity in the tropical zone of Africa is typically
70--90\% RH, exceeding 95\% RH during the rainy season.
Under such conditions, bees are not biologically capable of
reducing honey moisture to $\leq$$20\%$ solely through natural
hive ventilation~\cite{tadele2025}. The beekeeper faces a choice:
early harvesting and technological dehumidification, or the risk
of in situ fermentation or plundering by pests. Tadele et al.~\cite{tadele2025}
classify dehumidification as a \hl{technological necessity},
not an intent to adulterate, emphasizing that the alternative
is crop loss or fermented product. Studies show that honey
subjected to proper dehumidification retains full enzymatic
and antioxidant profile~\cite{lastriyanto2020,spel3upm}.

\subsubsection*{African Production: Climatic Humidity and Traditional Hive Systems}

In Ethiopia, 95.5\% of hives are traditional, with only 0.2\%
being frame hives~\cite{ethaau2007}. Paradoxically, the study
by Gobessa et al.~\cite{sciendo2012} demonstrated that honey
from traditional hives is characterized by \hl{significantly
lower} moisture than from modern hives ($p < 0.001$), as
harvesting occurs once or twice per year, giving more time
for natural evaporation. At the same time, traditional hives
make frame inspection and assessment of the degree of capping
impossible without destroying the combs; applying a frame-based morphological criterion to these systems illustrates the limited transferability of maturity indicators derived mainly from temperate-climate European production systems \cite{tadele2025}.


\subsection{Honeys of Stingless Tropical Bees}


Stingless bees (Meliponini, >500~species) produce honey with
moisture typically of 25--35\% and $a_w = 0.60$--$0.86$,
a different sugar profile, and higher acidity than
\textit{Apis mellifera} honeys~\cite{stinglessyeastbrazil2021,stinglesswjava2024}.
Meliponini do not possess nectar dehydration mechanisms
comparable to the European bee: smaller honey crops and less
developed hive ventilation are biological characteristics,
not production deficiencies. A study of 17 species
from Brazil~\cite{stinglessyeastbrazil2021} demonstrated a range
of $a_w = 0.60$--$0.86$, with osmophilic yeasts isolated
from these honeys characterized by high osmotolerance and low
fermentative capacity---biological adaptation to an environment
with elevated $a_w$. Meliponini honeys require separate
standardization; assessing their quality on the basis of
\textit{Apis mellifera} standards is biologically
unjustified~\cite{stinglessyeastbrazil2021,stinglesswjava2024}.


\subsection{Specific Microclimates and Weather Conditions in the Temperate Zone}

Honey extraction with moisture levels that pose a real risk of
fermentation also occurs in the temperate zone, including in
Europe. Many areas from which honey is harvested are characterized
by a specific microclimate with persistently elevated air humidity.
Weather factors such as rainfall lasting for periods of several
days to several weeks influence elevated water content in combs.
This is precisely why the Apimondia position states that a justified
beekeeping practice is to leave honey after extraction in open
vessels for the purpose of evaporating excess water. The extraction of
honey with elevated water content is therefore a fact, as is human
intervention aimed at removing excess water for the purpose of
microbiological stabilization. Commonly practiced by European
beekeepers is not only the periodic mixing of honey to increase the rate
of evaporation but also the use of various technical devices for
this purpose (e.g., rotating discs that increase the evaporation
surface). Technical development and process automation are visible
at all stages of honey extraction and serve not only to increase
efficiency but also effectively enhance product safety.
Unquestionably, however, they cannot be instruments of honey
adulteration. Therefore, clear, accessible, and verified criteria
contained in normative acts should underpin the assessment of such
activities, and within justified scope, prohibited practices should
be precisely specified (e.g., the use of nanofiltration, ion exchange
resins, and vacuum dehydration).

\subsection{Comparison of Production Systems}


Table~\ref{tab:hive_systems} summarizes the key characteristics of
the three globally dominant hive types. Only the Langstroth hive
provides full control over honey ripening; KTBH and traditional
hive systems are structurally limited in this respect, meaning that
honeys from these systems may fail to meet the capping criterion
while complying with all good beekeeping
practices~\cite{tadele2025,sciendo2012,lrrd2013}.

\begin{table}[H]
\caption{Comparison of beekeeping production systems with respect to
         honey maturity control. Source: own compilation
         based on~\cite{tadele2025,sciendo2012,lrrd2013}.}
\label{tab:hive_systems}
\begin{adjustwidth}{-\extralength}{0cm}

\begin{tabularx}{\fulllength}{ccCC}
\toprule
\textbf{Feature} &
\textbf{Langstroth Hive} &
\textbf{Kenya Top Bar Hive} &
\textbf{Traditional Hive} \\
\midrule
Capping control
    & Full (frame inspection)
    & Partial
    & None (destructive) \\
\midrule
In situ dehumidification
    & Yes
    & Limited
    & None \\
\midrule
Average yield [kg/hive/year]
    & 20--40
    & 15--26~\cite{lrrd2013}
    & 5--15~\cite{tadele2025} \\
\midrule
Typical honey moisture
    & 16--19\%
    & 17--20\%
    & 17--23\% \\
\midrule
Risk of yeast colonization
    & Low
    & Moderate
    & High \\
\midrule
Global use
    & Americas, Australia, Europe
    & East Africa
    & Africa, S.E. Asia \\
\bottomrule
\end{tabularx}
\end{adjustwidth}
\end{table}


\subsection{Limitations of the Morphological Capping Criterion}


Based on the above analysis,
{the criterion of full capping as the sole determinant of honey maturity has limited global applicability and may be technically inadequate in several production systems for five reasons:
} (1)~the climatic
impossibility of fully dehydrating nectar under conditions of
70--90\%~RH humidity~\cite{tadele2025}; (2)~the technological
inaccessibility of capping assessment in traditional hives without
comb destruction~\cite{sciendo2012}; (3)~the biological distinctness
of Meliponini honeys that precludes direct application of
\textit{Apis} standards~\cite{stinglessyeastbrazil2021};
(4)~the empirical inconsistency between the degree of capping and
the microbiological stability of honey~\cite{zhang2021,guo2020};
and (5)~the impossibility of gathering, archiving, and verifying
evidence, and of post factum proof from the degree of capping.
It is necessary to adopt an integrated model for assessing
functional maturity based on measurable parameters ($a_w$),
not on a morphological
criterion~\cite{tadele2025,codex_stan12,ruegg1981}.






\section{EU Honey Market and Regulatory Implications}

The European honey market is the world's largest regional import
market, and its structural characteristics---chronic domestic
production deficit, high dependence on imports from third countries,
and wide quality variation of supplied products---constitute
a context in which regulatory decisions concerning honey
disqualification criteria have direct economic, trade, and consumer
consequences. This section analyzes this market context as
justification for the precise approach advocated in this article,
based on functional maturity and water activity, in place of an
overly simplified approach based on a yeast count or morphological
criterion.


\subsection{Structure and Dynamics of Honey Imports into the EU}


The European Union is the world's largest regional honey importer,
generating a trade balance deficit in the beekeeping sector that
has no equivalent in any other region~\cite{cbi2024,pollinatorhub2025}.
EU production in the 27 Member States was estimated at approximately 272,286 tonnes in 2022–2023, covering only 55–60\% of domestic demand, estimated at approximately 480--500 thousand tonnes~\cite{pollinatorhub2025,ec_honey_autumn2024}.
The difference is filled by imports from third countries.


In 2023, extra-EU honey imports amounted to approximately 163,700 tonnes, while the broader volume of honey entering and circulating within the EU supply chain, including intra-EU trade, was estimated at approximately 300,000–360,000 tonnes~\cite{cbi2024}.

In 2024, total imports (extra-EU + intra-Community) amounted to
approximately 290 thousand tonnes (a decline of 3.1\% year-on-year
after three years of growth), but the long-term trend remains
upward: the EU market is expected to reach a volume of 445 thousand
tonnes by 2035 (CAGR +1.9\%)~\cite{eurostat_honey2023}.

The main external suppliers in 2023 were: \hl{China} (37\% of extra-EU imports), \hl{Ukraine} (28\%),
\hl{Argentina} (approx.\ 7\%), \hl{Mexico} and
\hl{Cuba}~\cite{tridge2024}.
The largest national importers within the EU were Germany
(41 thousand tonnes, 25\% of extra-EU imports), Belgium
(31.4 thousand tonnes, 19\%), Poland (23.3 thousand tonnes, 14\%),
Spain (15.7 thousand tonnes), and France
(7.7 thousand tonnes)~\cite{eurostat_honey2023}.


\subsection{The Scale of Honey Adulteration on the EU Market: Empirical Data}


The scale of the honey authenticity problem on the EU market is
documented by a number of studies conducted at the institutional
level. The coordinated control action of the European Commission
``From the Hives'' (2021--2022), covering analysis of 320 samples
of imported honey, demonstrated that \hl{46\%} of the analyzed
samples were suspected of non-compliance with the Honey Directive,
representing a threefold increase compared to a similar study from
2015 to 2017, when the figure was 14\%~\cite{eufromhives2021}.
The highest absolute percentage of suspect shipments concerned China
(74\%), while Turkey showed the highest relative rate
(93\%)~\cite{eufromhives2021}. Among the identified adulteration
practices, the addition of sugar syrups (from sugarcane, corn, and rice)
dominated, along with blending authentic honey with syrup products
in proportions below the detection threshold of conventional
\mbox{methods~\cite{eufromhives2021,schoder2026}.}

These data indicate a real and serious honey authenticity problem
on the EU market but simultaneously clearly profile the dominant
form of adulteration as \hl{sugar syrup adulteration}, not as
the distribution of honey with an elevated yeast count. JRC and OLAF
documents do not indicate a single case of honey disqualification
based solely on the yeast CFU/g count without simultaneous
determination of active fermentation or syrup
adulteration~\cite{eufromhives2021}. This is a significant
observation from the perspective of this article: adulteration
practices focus on the chemical dilution of honey, not on the manipulation
of its microflora.

In response to this crisis, the European Commission adopted
Directive (EU) 2024/1438 (part of the ``Breakfast Directives''
package), introducing from 2026 an obligation to declare all
countries of origin of honey blend components in the principal
field of view of the label~\cite{cbi_traceability2024}. The
implementation of this requirement strengthens the traceability
infrastructure but does not resolve the problem of the lack of
validated parameters specific to immaturity.


\subsection{Analytical Gaps: Absence of a Validated Immaturity Biomarker}


The identification of immature honey as a distinct category of
adulteration requires the existence of analytical tools capable of
reliably distinguishing between immature and mature honey, and here
science encounters a fundamental limitation. As the review by
Tsagkaris et al.~\cite{tsagkaris2021} demonstrated, standard
physicochemical methods (moisture, diastase, invertase, HMF,
and sugar profile) are not sufficient for unambiguous determination
of immaturity, because:

\begin{enumerate}
    \item \hl{There is no specific marker:}
          no single chemical compound or enzyme is identified
          as a biomarker exclusively of immaturity in a manner
          independent of botanical and geographical type;
          the physicochemical parameters of mature honey from
          the tropics may resemble the parameters of immature
          honey from the temperate zone~\cite{tsagkaris2021,guo2020}.
    \item \hl{Metagenomics and metabolomics have limited
          discriminatory power without reference databases:}
          as Schoder~\cite{schoder2026} emphasizes in his 2026
          review, the high analytical resolution of omic methods
          (genomics, proteomics, NMR, GC-MS, and LC-HRMS) does not
          guarantee unambiguous fraud attribution without solid,
          globally representative reference databases;
          \hl{``high analytical resolution alone therefore
          does not guarantee unambiguous fraud
          attribution''}~\cite{schoder2026}.
    \item \hl{Blending effects mask immaturity:}
          blending immature with mature honeys in proportions
          below 50\% is analytically indistinguishable from mature
          honey by conventional methods~\cite{schoder2026}; it
          should be emphasized that this problem also applies to
          the honey extraction stage---it is impossible to
          determine what fraction of honey was mature and what
          was immature at the time of extraction.
    \item \hl{Yeast presence is not a specific marker:}
          as demonstrated in Section~\ref{4}, the osmophilic yeast count
          in honey is determined mainly by climatic production
          conditions and hive systems, not by the degree of
          maturity, which excludes CFU/g as a useful immaturity
          biomarker~\cite{rodriguezmachado2024,schoder2026}.
\end{enumerate}

The most recent studies point to promising directions in the search
for immaturity biomarkers: Sun et al.~\cite{sun2021} identified
decenedioic acid and other fatty acids of bee origin as potential
molecular markers whose concentration correlates with the degree of
maturity in a manner independent of yeast count. NMR-based
metabolomics~\cite{nmrhoney2023} enables simultaneous quantification
of several dozen compounds in a single measurement, providing
a basis for developing a multidimensional maturity index. However,
to date, none of the proposed biomarkers has undergone full
multilaboratory validation encompassing the global diversity of
botanical and geographical types~\cite{tsagkaris2021}.

It should be borne in mind that legally established methods should
be as simple, rapid, and accessible as possible since they will
constitute the basic tool of every beekeeper (similar to the current
refractometer) for determining the timing of extraction, as it is
this stage that determines whether honey is mature or not.


\subsection{Regulatory Risk: Consequences of Overly Strict Interpretation}


Adopting a disqualification criterion based on yeast CFU/g count,
without taking into account water activity, climatic production
conditions, and the absence of active fermentation entails a number
of serious regulatory and market risks, whose analysis is necessary
for a proper assessment of the implications of the Apimondia
position:

\begin{itemize}
        \item {\hl{Risk} %MDPI: We formatted this as a list. Please confirm. The same applies to the following highlights.
%I confirm
 of discrimination against tropical producers.}
        \end{itemize}

As demonstrated in Section~\ref{5}, honeys from tropical and traditional
systems naturally contain higher concentrations of osmophilic yeasts
due to hive type, climatic conditions, and richer endemic
microbiological
biodiversity~\cite{tadele2025,stinglessyeastbrazil2021}.
Disqualification of these honeys based on CFU/g would effectively
result in excluding from the EU market products legally produced
in countries of East Africa, Southeast Asia, and Latin America,
while not affecting syrup-adulterated honeys (which contain a
standard or reduced yeast count following the ultrafiltration
process)~\cite{eufromhives2021}.

\begin{itemize}
        \item {\hl{Risk} 
 of aggravating the deficit and price increases.}
        \end{itemize}

Excluding from the EU market honeys from China, Ukraine, and
Argentina---the three largest suppliers collectively accounting
for $\approx72\%$ of extra-EU imports~\cite{tridge2024}---based
on the CFU/g criterion would cause a sharp rise in honey prices
on the European market and an aggravation of the deficit,
particularly severe for consumers in countries with low domestic
consumption (Poland, Germany, and Belgium as net
importers)~\cite{eurostat_honey2023}.

\begin{itemize}
        \item {\hl{Risk} 
 of creating arbitrary trade barriers.}
\end{itemize}

Applying a criterion that has no basis in Codex Alimentarius or
EU Directive 2001/110/EC can be challenged by trading partners
under the WTO-SPS Agreement (Agreement on the Application of
Sanitary and Phytosanitary Measures), which requires that
trade-restrictive measures have a scientific basis and be
proportionate to the identified risk~\cite{codex_stan12,eudir2001110}.
The risk of harvesting immature honey exists for all honeys,
regardless of where the apiary is located. Introduced requirements
and restrictions should have a global dimension and their
proportional reflection in EU law with respect to EU honey producers.


\begin{itemize}
        \item {\hl{Risk} 
 of missing the actual problem.}
\end{itemize}

As demonstrated by the ``From the Hives'' action, the dominant
form of honey adulteration in the EU is sugar syrup adulteration,
not the distribution of immature honey with an elevated yeast
count~\cite{eufromhives2021}.
Focusing inspection resources on CFU/g measurement at the expense
of isotope ratio analysis ($\delta^{13}$C, $\delta^{2}$H,
$\delta^{18}$O) and NMR fingerprinting would lead to the allocation
of inspection resources inconsistent with the actual adulteration
risk profile.
It should be noted, however, that the European Commission later
withdrew from treating $^1$H-NMR profiling as a sufficiently
established official control method, emphasizing the need for
further research, validation, and harmonization of analytical
methods~\cite{EuropeanCommission2023NMR}.



\subsection{Regulatory Proposals: Harmonization, Traceability, and Proportionality}


Based on the above market and regulatory analysis, this article
formulates the following proposals addressed to standardization
bodies (ISO, Codex), EU institutions (DG SANTE, EC), and industry
organizations (Apimondia, Copa-Cogeca):

\begin{enumerate}
    \item \highlighting{Update EU Directive 2001/110/EC with a water
          activity (\boldmath{$a_w$}) criterion:}
          introduction of mandatory $a_w$ measurement alongside
          moisture content as a primary parameter for fermentation
          risk assessment. Measurement of $a_w$ is a validated,
          commercially available method (instruments such as
          AquaLab, Rotronic HygroLab) and inexpensive; this would
          constitute a substitution of the morphological criterion
          with a physicochemical criterion having direct scientific
          justification~\cite{ruegg1981,zamora2006}.

    \item \hl{Establish 
 a multiparametric microbiological
          assessment protocol:} fermentation risk assessment should
          be conducted simultaneously for $a_w$, water content,
          storage temperature, and yeast count (in accordance with
          the risk Equation (\ref{eq:fermentation_risk}) proposed
          in Section~\ref{4.3}); none of these parameters may be
          decisive in isolation from the
          others~\cite{analytica2020,apinz2021}.


\end{enumerate}






\section{Proposed Framework for Functional Maturity Assessment}

The preceding sections indicate that honey maturity is a continuous and multidimensional phenomenon. Therefore, no single parameter, whether morphological, microbiological, or physicochemical, is sufficient for its reliable assessment. Rather than treating honey as simply “mature” or “immature”, it is more appropriate to assess the degree of functional maturity in relation to microbiological stability, chemical quality, and fermentation risk. Natural variability related to botanical origin, climate, weather conditions, hive type, and beekeeping practices makes rigid universal classification difficult. However, the main quality and commercial risk associated with insufficient maturity, namely fermentation, can be assessed using measurable parameters, particularly water activity. The framework proposed in this section should be regarded as a conceptual and operational tool requiring further validation before possible use in official control practice.



\subsection{Water Activity as a Priority Parameter}

Among all the biomarkers mentioned, water activity
($a_w$) is proposed as the priority parameter for three reasons.
First, it is a thermodynamic measure of the \hl{availability
of water to microorganisms}, not merely the water content per unit
mass of product; it directly correlates with the biological
possibility of osmophilic yeast growth, making it the only
parameter that can replace the entire range of indirect
microbiological indicators~\cite{ruegg1981,waaw2006}.
Second, measurement of $a_w$ is rapid (3--5 min), inexpensive,
and can be performed at an apiary or control laboratory using a
standard instrument such as AquaLab CX3TE or Rotronic HygroLab;
the result is reproducible and unambiguous, requiring no expert
interpretation~\cite{zamora2006}.
Third, the criterion $a_w < 0.60$ is based on a biologically
justified mechanism---it is the physiological threshold of
osmophilic yeasts determined in vitro under isobaric
conditions~\cite{waaw2006,ruegg1981}---and not on an arbitrarily
chosen administrative value.

It is nevertheless necessary to conduct studies on whether,
similarly to permitted exceptions regarding water content, exceptions
should also be established in this regard. The microbiological
stability of certain honey types may result not so much from low
water activity as from specific components present in plant nectar.


\subsection{Hierarchical Decision Framework}

{
Based on the biomarkers described above, the proposed functional maturity assessment framework can be structured into two main analytical levels, organized according to the principle of proportionality. Level I provides basic screening based on moisture and water activity, whereas Level II provides an extended qualitative assessment based on enzymatic and chemical parameters for borderline or uncertain cases. Candidate metabolomic biomarkers are treated as a research-oriented validation layer requiring further multilaboratory verification rather than as routine official control parameters.
}

\subsubsection{Level I ---Basic Screening}

Level I constitutes a first-order filter based on two parameters
measured directly on the product:

\begin{itemize}
    \item \highlighting{Moisture \boldmath{$\leq 20\%$}} taking into account
          legally specified exceptions (refractometry,
          Codex STAN 12-1981 method), a normative parameter,
          mandatory in all major standards;
    \item \highlighting{Water activity \boldmath{$a_w < 0.60$}} taking into
          account possible exceptions, a parameter proposed as
          a normative supplement; when both of these criteria
          are met, the honey is \hl{microbiologically stable
          by virtue of the state of microbiological knowledge},
          and yeast growth is  {not observed and is effectively inhibited under standard storage conditions,}
          regardless of their count~\cite{ruegg1981,waaw2006}.
\end{itemize}

\hl{Key rule of Level I:} if moisture $\leq$ $20\%$
\textsc{AND} $a_w < 0.60$, the product is not subject to
disqualification based on yeast count (CFU/g), regardless of the
determined value of this parameter. Disqualification based solely
on CFU/g, without confirmation of exceeding the $a_w$ threshold,
is scientifically unjustified and legally
inadequate~\cite{codex_stan12,eudir2001110}.

If Level I indicates a borderline state ($18.1$--$20.0\%$
moisture or $a_w = 0.60$--$0.62$) or the yeast count exceeds
$1000$~CFU/g, analysis is escalated to Level~II.

\subsubsection{Level II ---Extended Qualitative}

{The hierarchical decision structure of the proposed functional maturity assessment framework, including the transition from basic screening to extended qualitative assessment, is summarized in Figure \ref{fig:fmam}}.
Level II provides an assessment of functional maturity in all
three dimensions defined in Section~\ref{3.3}: microbiological, chemical,
and biological stability. Level~II parameters include:

\begin{enumerate}
    \item {\hl{Enzymatic profile:}}
    \begin{itemize}
        \item Diastase activity $\geq$ $8$~Schade units (or $\geq$$3$~units
              for botanical types with naturally low amylolytic
              activity, e.g., acacia, citrus, lavender)~\cite{codex_stan12}.
    \end{itemize}
    \item {\hl{Chemical parameters:}}
    \begin{itemize}
        \item F/G ratio $\geq$ $0.9$; HMF $\leq$ $40$~mg/kg;
              free acidity $\leq$ $50$~mEq/kg
              (Codex)~\cite{codex_stan12}.
    \end{itemize}
\end{enumerate}


\begin{figure}[H]
\centering
\tikzset{
  fmbox/.style={rectangle, rounded corners=4pt, draw=black, fill=gray!10,
                text width=4.2cm, align=center, minimum height=0.7cm, font=\footnotesize},
  levelbox/.style={rectangle, rounded corners=4pt, draw=blue!70!black, fill=blue!8,
                   text width=5.0cm, align=center, minimum height=1.0cm, font=\footnotesize},
  okbox/.style={rectangle, rounded corners=4pt, draw=green!60!black, fill=green!8,
                text width=3.4cm, align=center, minimum height=0.7cm, font=\footnotesize},
  warnbox/.style={rectangle, rounded corners=4pt, draw=orange!80!black, fill=orange!8,
                  text width=3.4cm, align=center, minimum height=0.7cm, font=\footnotesize},
  fmarrow/.style={-{Stealth[length=2mm]}, thick}
}
\resizebox{\linewidth}{!}{%
\begin{tikzpicture}[node distance=0.6cm and 0.8cm]
  \node[fmbox] (start)
    {Honey sample};

  \node[levelbox, below=of start] (lI)
    {\textbf{Level I -- Basic Screening}\\
     Moisture $\leq 20\%$ \textsc{and} $a_w < 0.60$\\
     \scriptsize(refractometry + $a_w$ meter)};

  \node[okbox,  below left =1.2cm and 0.6cm of lI] (pass1)
    {Microbiologically stable\\No CFU/g disqualification};

  \node[warnbox, below right=1.2cm and 0.6cm of lI] (border)
    {Borderline or CFU/g\,$> 1000$\\$\Rightarrow$ escalate to Level~II};

  \node[levelbox, below=1.0cm of border] (lII)
    {\textbf{Level II -- Extended Qualitative Assessment}\\
     Diastase $\geq 8$\,Schade u.;\; HMF $\leq 40$\,mg/kg\\
     F/G $\geq 0.9$;\; Acidity $\leq 50$\,mEq/kg\\
     Metabolomic indicators: validation context
     };

  \node[okbox,  below left =1.2cm and 0.5cm of lII] (pass2)
    {Functionally mature\\Compliant};

  \node[warnbox, below right=1.2cm and 0.5cm of lII] (fail2)
    {Parameter(s) exceeded\\Further investigation};

  \draw[fmarrow] (start) -- (lI);
  \draw[fmarrow] (lI) -- node[above left,  font=\scriptsize]{Both met}   (pass1);
  \draw[fmarrow] (lI) -- node[above right, font=\scriptsize]{Borderline / fail} (border);
  \draw[fmarrow] (border) -- (lII);
  \draw[fmarrow] (lII) -- node[above left,  font=\scriptsize]{All met}       (pass2);
  \draw[fmarrow] (lII) -- node[above right, font=\scriptsize]{Any exceeded}  (fail2);
\end{tikzpicture}}
\caption{{\hl{Hierarchical} %MDPI: 1. Please note that changes to the position/size of figures or tables may occur during the production stage. 2. Please note that figures and captions require careful review at this stage as well. Please verify them thoroughly before returning your proofread manuscript. Doing so helps us avoid an additional round of revisions and ensures timely publication.
 decision structure of the proposed functional maturity assessment framework (FMAM). Level I provides rapid basic screening based on moisture content and water activity $(a_w)$. Level II provides extended qualitative assessment for borderline or analytically uncertain cases using enzymatic and chemical parameters. Candidate metabolomic biomarkers are considered a~research-oriented validation component requiring further multilaboratory verification. The framework is a conceptual model and should not be regarded as a validated official control procedure at this stage.}}
 %Figure and caption verified. All parameter values are consistent
%with the text and Codex STAN 12-1981. One correction made in the
%Level II node: "SU" changed to "Schade u." to match the
%abbreviation used in Table 1 and the main text.

\label{fig:fmam}
\end{figure}


\subsection{Recommendations for Regulatory Bodies and Apimondia}



Based on the proposed FMAM and the conducted scientific analysis,
this article formulates the following recommendations addressed
to key stakeholders:


\subsubsection{Recommendations for Apimondia}

\begin{enumerate}
    \item {\highlighting{Supplement the 2023 position with the \boldmath{$a_w$}
          criterion:}} the Apimondia recommendation should
          explicitly state that the presence of osmophilic yeasts
          in honey does not constitute grounds for declaring
          immaturity, provided moisture $\leq$ $20\%$ and
          $a_w < 0.60$; both parameters should be listed as
          simultaneous necessary
          conditions~\cite{ruegg1981,waaw2006}.
    \item {\hl{Take into account the regional diversity of
          production systems:}} the position should contain a clear
          clause regarding honeys from tropical climates and
          traditional systems, recognizing technological
          dehumidification as a legal production practice~\cite{tadele2025}
          or excluding specific dehumidification techniques
          that negatively affect the product.
    \item {\hl{Initiate dialogue with ISO TC~34/SC~17:}} %MDPI: 1. The standard should be listed in the reference part and cited when it appears in the main text (in the form of a reference number, e.g., [2]). Please revise. 2. Please provide detailed information here, and we can help you add it to the reference part and rearrange the citation order. e.g., TB 10099-2017; Code for Design of Railway Station and Terminal. China Railway Publishing House: Beijing, China, 2017.
    
    %Thank you for the comment. We would like to clarify that ISO TC 34/SC 17 is not a standard but an ISO technical subcommittee (Management Systems for Food Safety). It does not require a bibliographic reference. However, we have added a reference to the committee's webpage for transparency: ISO/TC 34/SC 17. Management Systems for Food Safety. International Organization for Standardization. Available online: https://www.iso.org/committee/583916.html
          the ISO 24607:2025 standard on bee honey should be
          extended to include the $a_w$ criterion as a normative
          parameter~\cite{codex_stan12}.
\end{enumerate}

\subsubsection{Recommendations for EU Bodies (EC, DG SANTE)}


{\hl{Consider the inclusion of} \highlighting{\boldmath{$a_w$}} \hl{as a risk-based control parameter in official honey quality controls}}, particularly for consignments with borderline moisture values, products from high-humidity production systems, or cases in which fermentation risk is analytically uncertain.


\subsubsection{Recommendations for Laboratories and the Scientific Community}

\begin{enumerate}
    \item {\hl{Standardize the multiparametric microbiological
          assessment protocol}} encompassing the simultaneous measurement
          of CFU/g, $a_w$, temperature, and phase state of honey
          (liquid/crystallized); results should be reported
          jointly, not as isolated values~\cite{zamora2006,analytica2020}.
    \item {\hl{Perform multilaboratory validation of the
          decenedioic acid biomarker}} and extend its verification
          beyond three botanical types (rapeseed, acacia, and jujube)
          to tropical, stingless bee, and traditional production
          system honeys~\cite{sun2021}.
    \item {\hl{Develop portable in situ analysis tools:}}
          miniaturized NIR and Raman spectrometers allow
          measurement of moisture, $a_w$, and sugar profile
          directly at the apiary or at the border, with accuracy
          comparable to laboratory
          methods~\cite{honeyid2025,schoder2026}.
\end{enumerate}



 


\section{Conclusions and Future Perspectives}

The conclusions of this review are presented in four complementary parts, covering the scientific interpretation of honey maturity, the regulatory implications of the Apimondia position, the proposed functional maturity assessment framework, and the limitations requiring further research.

This article undertook an analysis of the Apimondia position
on the production and distribution of immature honey, with
particular attention to the microbiological, regulatory, and
market implications of using osmophilic yeast count and the
morphological criterion of comb cell capping as quality indicators.
The conducted analysis has led to a number of conclusions of varying
degrees of generalization, from specific biochemical findings
to normative recommendations of global scope.


\subsection{Main Scientific Conclusions}


{\hl{(C1) The presence of osmophilic yeasts in honey is a biologically normal phenomenon.}}
Yeasts of the genera \textit{Zygosaccharomyces}, \textit{Schizosaccharomyces},
and \textit{Metschnikowia} are an inseparable component of the hive
ecosystem, regardless of production method, hive system, and country
of origin. The natural yeast population count in correct, non-fermenting
commercial honey may amount to several hundred to a few thousand CFU/g
without any effect on the microbiological stability of the product,
provided water activity $a_w < 0.60$~\cite{ruegg1981}.
This conclusion has been well established in the literature for decades
and confirmed by all current systematic
reviews~\cite{biochemreactions2022,rodriguezmachado2024}.

{\hl{(C2) Water activity, not moisture content or yeast count, is the primary fermentation risk parameter.}}
The physiological growth threshold of osmophilic yeasts
($a_w = 0.61$--$0.62$) defines a biologically justified
boundary below which honey fermentation is  {thermodynamically unfeasible and has not been
observed under standard storage and temperature conditions,}
regardless of the yeast count expressed in
CFU/g~\cite{waaw2006,ruegg1981}.
Glucose crystallization constitutes an additional factor that
dynamically modifies the $a_w$ of the liquid phase, meaning that
fermentation risk assessment must take into account the phase state
of the product and cannot be based solely on the global water
content~\cite{zamora2006,ijfans2019}.

{\hl{(C3) Honey ripening is a continuous biological process, not a binary event.}}
Maturity indicators---invertase, glucose oxidase, and diastase
activity, organic acid profile, and fructose-to-glucose ratio---change gradually during the ripening process and do not correlate
unambiguously with the moment of cell capping by
bees~\cite{zhang2021,sun2021}.
Uncapped honey may exhibit a chemical profile appropriate for
mature honey if moisture is low; capped honey may lose maturity
characteristics as a result of inappropriate heat
treatment~\cite{codex_stan12,guo2020,zhang2021}.

{\hl{(C4) Some honey maturity criteria are derived mainly from temperate-climate, frame-hive production systems and may have limited applicability under diverse global production and microclimatic conditions.}}
The climatic, technological, and biological conditions of honey
production in tropical and subtropical zones, including high
environmental humidity, log and basket hive systems, and the
distinct properties of stingless bee honeys (Meliponini), mean
that uniform morphological (capping) and microbiological (CFU/g
without $a_w$) criteria cannot serve as global quality indicators
without systematically favoring honeys produced in European
systems~\cite{tadele2025,stinglessyeastbrazil2021,sciendo2012}.
This position is consistent with the general trend in food
legislation toward recognizing the diversity of local production
systems and moving away from norms derived exclusively from
Western European practices~\cite{phiri2022}.

{\hl{(C5) The dominant form of honey adulteration on the EU market is sugar syrup adulteration, not the distribution of honey with elevated yeast count.}}
Data from the coordinated control action of the European Commission
``From the Hives'' (2021--2022) unambiguously indicate that 46\%
of suspect imported honey shipments contained additions of
C$_3$/C$_4$ syrups or oligosaccharides~\cite{eufromhives2021}.
None of the available control reports identifies elevated CFU/g
count as the primary indicator of adulteration in inspection
practice. Allocation of inspection resources toward $a_w$
measurement is therefore more appropriate to the actual market
risk profile than the intensification of microbiological
methods~\cite{schoder2026,eufromhives2021}.


\subsection{Regulatory and Practical Implications of the Apimondia Position}


The Apimondia position of 2023 on immature honey is a valuable and
necessary contribution to the protection of the integrity of the
beekeeping market. The threat posed by systemically immature honey,
produced by the method of industrial nectar dehydration in a short
time and without maintaining the biological ripening processes, is
real and constitutes unfair competition with producers adhering to
quality standards. In this respect, the Apimondia position addresses an important and legitimate concern related to honey integrity and fair trade.

This article does not undermine this intention but indicates that
the current formulation of the position is analytically incomplete
in two important areas. First, it does not distinguish between
(a)~fermentation as a product characteristic at the time of harvest,
and (b)~fermentation as a post-harvest event resulting from supply
chain errors. Second, it does not take into account the fundamental
role of water activity as a parameter superior to yeast count in the
assessment of fermentation risk. Supplementing the Apimondia position
with these two elements will not weaken its protective message but
will instead strengthen it scientifically and make it more resistant
to legal challenges in trade disputes with importers from tropical
countries~\cite{ruegg1981,tadele2025,analytica2020}.

{Furthermore, an overly broad interpretation of yeast presence as an indicator of immaturity may lead to disproportionate assessment outcomes for honeys in which naturally elevated osmophilic yeast counts are frequent and not associated with active fermentation.}


\subsection{Functional Maturity Assessment Framework}


The proposed functional maturity assessment framework (FMAM) offers a conceptual and operational basis for further development of honey quality assessment because it is:

\begin{itemize}
    \item {\hl{Scientifically justified}}, based on the
          thermodynamic properties of water activity, enzymatic
          ripening kinetics, and
          metabolomics~\cite{ruegg1981,zhang2021,guo2020};
    \item {\hl{Compliant with applicable regulations}}---Codex STAN 12-1981 and EU Directive 2001/110/EC,
          and does not require their revision, only interpretive
          extension~\cite{codex_stan12,eudir2001110};
    \item {\hl{Technically accessible}}, as $a_w$ measurement
          and refractometry can be performed in any control
          laboratory without specialized
          infrastructure~\cite{waaw2006,zamora2006};
    \item {\hl{Potentially globally applicable}}, since criteria based
          on $a_w$ and enzymatic profile are independent of the
          hive system and climatic conditions of the place of
          \mbox{production~\cite{tadele2025,stinglessyeastbrazil2021};}
    \item {\hl{Proportionate to the risk}}, preventing
          escalation to costly
          methods~\cite{schoder2026,nmrhoney2023}.
\end{itemize}


\subsection{Limitations and Directions for Future Research}


The authors of the article are aware of the limitations of this
study. The proposed FMAM model is based on empirical data available
at the time of manuscript preparation and requires further
verification through prospective multilaboratory studies.
The following directions for further research are particularly
indicated:

\begin{enumerate}
    \item {\hl{Multilaboratory validation of the decenedioic
          acid biomarker}} on samples representing the global
          diversity of botanical types and production systems,
          extending beyond the three types (rapeseed, acacia,
          and jujube) studied by Sun et al.~\cite{sun2021}; validation
          should encompass stingless bee honeys (Meliponini)
          and honeys from traditional East African
          systems~\cite{tadele2025}.

    \item {\hl{Development and validation of a multidimensional
          functional maturity index}} (FMI,
          {\hl{Functional Maturity Index}}) integrating the
          results of Level~II parameters into a single index
          with a value of 0--100, calibrated for different
          botanical and geographical types; such an index could
          replace binary disqualification decisions with a
          gradient approach proportionate to the deviation from
          the optimal state~\cite{zhang2021,guo2020,sun2021}.

    \item {\hl{Harmonization of ISO TC~34/SC~17 standards
          taking into account the meliponiculture production
          system}}---the stingless bee honey category requires
          a separate standardization track with its own moisture
          thresholds, $a_w$, and enzymatic criteria, taking
          into account the biological properties of these
          insects~\cite{stinglessyeastbrazil2021,stinglesswjava2024}.
\end{enumerate}

Undertaking this research is possible only within the framework
of international cooperation encompassing scientific institutions,
beekeeping organizations, and regulatory bodies from countries of
both the ``Global North'' and tropical countries that are the
main honey producers. This dialogue is not only a scientific need but a necessary condition for building a globally equitable and
scientifically credible quality assessment system for one of the
oldest food products manufactured by humankind.


\vspace{12pt}
 
\authorcontributions{Conceptualization, A.G., B.L., P.R., and J.K.B.
methodology, A.G., B.L., and P.R.
investigation, A.G., B.L., and P.R.,
resources, A.G., B.L., P.R.,
formal analysis, A.G., B.L., P.R., 
writing---original draft preparation, A.G., B.L., P.R., J.K.B., R.P.-F., and E.T.
writing---review and editing, A.G., B.L., P.R., J.K.B. 
supervision, A.G., B.L., P.R., J.K.B., R.P.-F.
project administration, B.L., P.R., J.K.B. 
funding acquisition, P.R.
All authors have read and agreed to the published version of the manuscript.}

\funding{\hl{No} %MDPI: Information regarding the funder and the funding number should be provided. Please check the accuracy of funding data and any other information carefully.
 external funding was received for the preparation of this manuscript. The article processing charge (APC) will be covered by by Stowarzyszenie Polska Izba Miodu (Polish Honey Chamber Association). The APC funder had no role in the preparation of the manuscript or in the decision to submit it for publication.}
%Thank you for the comment. We confirm that no external funding was received for the research presented in this manuscript. The article processing charge (APC) is covered by Stowarzyszenie Polska Izba Miodu (Polish Honey Chamber Association). The APC funder had no role in the design of the study, collection, analysis, or interpretation of data, preparation of the manuscript, or in the decision to submit it for publication.






\institutionalreview{Not applicable.}
\informedconsent{Not applicable.}
\dataavailability{No new data were created or analyzed in this study. Data sharing is not applicable to this article.}

\acknowledgments{\hl{During} %MDPI: Please ensure that all individuals included in this section have consented to the acknowledgement.
 the preparation of this work, the authors used Claude Sonnet 4.6 (Anthropic) in order to improve the readability and language of the manuscript. The authors have reviewed and edited the output and take full responsibility for the content of this publication.}

 %Thank you for the comment. We confirm that all individuals acknowledged in this section have consented to being acknowledged. We fixed a typo: Clause > Claude
\conflictsofinterest{The authors declare no conflicts of interest.}

%%%%%%%%%%%%%%%%%%%%%%%%%%%%%%%%%%%%%%%%%%
\begin{adjustwidth}{-\extralength}{0cm}

\reftitle{\highlighting{References} %MDPI: Please do NOT change the reference format. Our production editor has done thoroughly layout work for the reference. If any changes are necessary, please use the "highlight". Thanks for your cooperation.
}

\begin{thebibliography}{999}

\bibitem[{Apimondia Working Group on Adulteration of Bee
  Products}(2023)]{apimondia_fraud_v2}
{Apimondia Working Group on Adulteration of Bee Products}. 
\newblock {\emph{Apimondia Statement on Honey Fraud}; Version 2}; Apimondia Federation: \hl{Rome, Italy,} %MDPI: Newly added information. Please confirm. The following highlights are the same.
\newblock   2023.  
%I confirm

\bibitem[{CBI---Centre for the Promotion of Imports}(2024)]{cbi2024}
{CBI---Centre for the Promotion of Imports}.
\newblock {\emph{The European Market Potential for Certified Honey}}; 
\newblock CBI Market Intelligence: \hl{The Hague, The Netherlands,} 
  2024.
%Fixed - Rotterdam is incorrect, correct:  The Hague, The Netherlands

\bibitem[{EU Pollinator Hub}(2025)]{pollinatorhub2025}
{EU Pollinator Hub}.
\newblock {\emph{European Honey Market}};
\newblock EU Pollinator Hub: \hl{Brussels, Belgium,}  2025.
%I confirm

\bibitem[{Apimondia International Federation of Beekeepers'
  Associations}(2026)]{apimondia_web}
{Apimondia International Federation of Beekeepers' Associations}.
\newblock Apimondia Official Website.  2026.  Available online:  \url{https://apimondia.org/}  (accessed on 3 April 2026).

\bibitem[Garc{\'i}a(2025)]{garcia2025}
Garc{\'i}a, N.
\newblock {Apimondia Statement on Immature Honey Production}.
\newblock {\em Bee World} {\bf 2025}, {\em 102},~56--60.
\newblock {\url{https://doi.org/10.1080/0005772X.2025.2557068}}.

\bibitem[Munitis et~al.(1976)]{munitis1976}
Munitis, M.T.;  Cabrera, E.; Rodriguez-Navarro, A.
\newblock {An Obligate Osmophilic Yeast from Honey}.
\newblock {\em Appl. Microbiol.} {\bf 1976}, \hl{\emph{32}, 3.}
 \url{https://doi.org/10.1128/aem.32.3.320-323.1976}
%I confirm

\bibitem[{Codex Alimentarius Commission}(2023)]{codex_stan12}
{\emph{CODEX STAN 12-1981}}; 
\newblock {Codex Standard for Honey, Revised 2023}. Codex Alimentarius Commission; 
\newblock FAO/WHO: \hl{Geneva, Switzerland,}  2023.
%I confirm

\bibitem[Rodr{\'i}guez~Machado et~al.(2024)]{rodriguezmachado2024}
Rodr{\'i}guez~Machado, A.;  Caro, C.M.; Hurtado-Murillo, J.J.; Gomes, Lobo, C.J.; Zúñiga, R.N.; Franco, W.
\newblock {Unconventional Yeasts Isolated from Chilean Honey}.
\newblock {\em Foods} {\bf 2024}, {\em 13},~1582.
\newblock {\url{https://doi.org/10.3390/foods13101582}}.

\bibitem[Zamora and Chirife(2006)]{zamora2006}
\hl{Zamora, M.C.; Chirife, J.} %MDPI: Refs. 9 and 21 are duplicated. Please just confirm that all references are cited in main text and we will help you to remove duplicated references and rearrange the order.
%References duplicated, please fix it.
\newblock {Determination of Water Activity Change Due to Crystallization in
  Honeys from Argentina}.
\newblock {\em Food Control} {\bf 2006}, {\em 17},~\hl{59--64.}
\newblock {\url{https://doi.org/10.1016/j.foodcont.2004.09.003}}.

\bibitem[Rüegg and Blanc(1981)]{ruegg1981}
Rüegg, M.; Blanc, B.
\newblock The water activity of honey and related sugar solutions.
\newblock {\em Lebensm.-Wiss. Technol.} {\bf 1981}, {\em
  14},~1--6.

\bibitem[Zhang et~al.(2021)Zhang, Tian, Zhang, Li, Zheng, and Hu]{zhang2021}
Zhang, G.Z.; Tian, J.; Zhang, Y.Z.; Li, S.S.; Zheng, H.Q.; Hu, F.L.
\newblock Investigation of the Maturity Evaluation Indicator of Honey in
  Natural Ripening Process: The Case of Rape Honey.
\newblock {\em Foods} {\bf 2021}, {\em 10},~2882.
\newblock  {\url{https://doi.org/10.3390/foods10112882}}.

\bibitem[Guo et~al.(2020)]{guo2020}
Guo, N.;  Zhao, L.; Zhao, Y.; Li, Q.; Xue, X.; Wu, L.; Gomez Escalada, M.; Wang, K.; Peng, W.
\newblock {Comparison of the Chemical Composition and Biological Activity of
  Mature and Immature Honey: An HPLC/QTOF/MS-Based Metabolomic Approach}.
\newblock {\em J. Agric. Food Chem.} {\bf 2020}, {\em
  68},~4062--4071.
\newblock {\url{https://doi.org/10.1021/acs.jafc.9b07604}}.
%I confirm

\bibitem[Sun et~al.(2021)]{sun2021}
Sun, J.;  Zhao, H.; Wu, F.; Zhu, M.; Zhang, Y.; Cheng, N.; Xue, X.; Wu, L.; Cao, W.
\newblock {Molecular Mechanism of Mature Honey Formation by GC-MS- and
  LC-MS-Based Metabolomics}.
\newblock {\em J. Agric. Food Chem.} {\bf 2021}, \hl{\emph{69}, 3362--3370.} 
\newblock {\url{https://doi.org/10.1021/acs.jafc.1c00318}}.
%I confirm

\bibitem[{European Parliament and Council}(2001)]{eudir2001110}
{European Parliament and Council}.
\newblock {Council Directive 2001/110/EC of 20~December 2001 relating to
  honey}.
\newblock \emph{Off. J. Eur. Communities}  \textbf{2001},  \emph{L10}, 47--52,  .

\bibitem[Tsagkaris et~al.(2021)]{tsagkaris2021}
Tsagkaris, A.S.;  Koulis, G.A.; Danezis, G.P.; Martakos, I.; Dasenaki, M.; Georgiou, C.A.; Thomaidis, N.S.
\newblock {Honey Authenticity: Analytical Techniques, State of the Art and
  Challenges}.
\newblock {\em RSC Adv.} {\bf 2021}, {\em 11},~11273--11294.
\newblock {\url{https://doi.org/10.1039/D1RA00069A}}.

\bibitem[Schoder(2026)]{schoder2026}
Schoder, D.
\newblock Honey Fraud as a Moving Analytical Target: Omics-Informed
  Authentication Within a Multi-Layer Analytical Framework.
\newblock {\em Foods} {\bf 2026}, {\em 15},~712.
\newblock {\url{https://doi.org/10.3390/foods15040712}}.

\bibitem[Alaerjani et~al.()Alaerjani, Abu-Melha, Alshareef, Al-Farhan, Ghramh,
  Al-Shehri, Bajaber, Khan, Alrooqi, Modawe, and
  Mohammed]{biochemreactions2022}
Alaerjani, W.M.A.; Abu-Melha, S.; Alshareef, R.M.H.; Al-Farhan, B.S.; Ghramh,
  H.A.; Al-Shehri, B.M.A.; Bajaber, M.A.; Khan, K.A.; Alrooqi, M.M.; Modawe,
  G.A.;  et~al.
\newblock Biochemical Reactions and Their Biological Contributions in Honey.
\newblock {\em Molecules} \hl{\textbf{2022}, \emph{27}, 4719.} 
\newblock {\url{https://doi.org/10.3390/molecules27154719}}.
%I confirm

\bibitem[Balasubramanyam(2023)]{enzymesrole2023}
Balasubramanyam, M.V.
\newblock {Role of Different Enzymes in Nectar to Honey Transformations in
  Indigenous Rock Bee \textit{Apis dorsata} F.} In {\em Cutting Edge Research
  in Biology}; Book Publisher International: \hl{Bhanjipur, India,}  2023; Volume~6, pp. 130--141.
\newblock {\url{https://doi.org/10.9734/bpi/cerb/v6/17763D}}.
%Thank you for the comment. The publisher is Book Publisher International (BP International); its registered office is in London, UK (with an Indian office in Hooghly, West Bengal). We suggest listing the location as London, UK. We also note this chapter has a co-author (Misbhauddin Khan). Full details:
%Balasubramanyam, M.V.; Khan, M. Role of Different Enzymes in Nectar to Honey Transformations in Indigenous Rockbee, Apis dorsata F. In Cutting Edge Research in Biology Vol. 6; B P International: London, UK, 2023; pp. 130–141. https://doi.org/10.9734/bpi/cerb/v6/17763D


\bibitem[Tadele et~al.(2025)]{tadele2025}
Tadele, E.;  Worku, D.; Muluneh, T.; Ayana, Y.; Melese, A.
\newblock {Comprehensive Review on Improved Honey Production: Techniques,
  Challenges, Opportunities, and Future Prospects in Africa}.
\newblock {\em Front. Bee Sci.} {\bf 2025}, {\em 3},~1588416.
\newblock {\url{https://doi.org/10.3389/frbee.2025.1588416}}.

\bibitem[Biswas et~al.(2023)Biswas, Naresh, Jaygadkar, and
  Chaudhari]{nmrhoney2023}
Biswas, A.; Naresh, K.S.; Jaygadkar, S.S.; Chaudhari, S.R.
\newblock Enabling honey quality and authenticity with NMR and LC-IRMS based
  platform.
\newblock {\em Food Chem.} {\bf 2023}, {\em 418},~135825.
\newblock {\url{https://doi.org/10.1016/j.foodchem.2023.135825}}.

\bibitem[Zamora and Chirife(2006)]{waaw2006}
\hl{Zamora, M.C.; Chirife, J.}
\newblock Determination of water activity change due to crystallization in
  honeys from {Argentina}.
\newblock {\em Food Control} {\bf 2006}, {\em 17},~59--64.
\newblock {\url{https://doi.org/10.1016/j.foodcont.2004.09.003}}.
%I confirm

\bibitem[Xue et~al.(2025)Xue, Zhou, Guo, Zhao, Jiang, Huang, and
  Zhang]{yeastbiodiversity2025}
Xue, S.J.; Zhou, M.; Guo, J.; Zhao, F.Y.; Jiang, W.W.; Huang, X.; Zhang, J.Y.
\newblock Unveiling the culturable and non-culturable yeast biodiversity in
  chaste honey produced by the sympatric Apis cerana and Apis mellifera from
  eastern China.
\newblock {\em Int. J. Food Microbiol.} {\bf 2025}, {\em
  410},~111423.
\newblock {\url{https://doi.org/10.1016/j.ijfoodmicro.2025.111423}}.

\bibitem[Al-Waili et~al.(2012)Al-Waili, Salom, Al-Ghamdi, and
  Ansari]{alwaili2012}
Al-Waili, N.S.; Salom, K.; Al-Ghamdi, A.; Ansari, M.J.
\newblock Antibiotic, Pesticide, and Microbial Contaminants of Honey: Human
  Health Hazards.
\newblock {\em  Sci. World J.} {\bf 2012}, {\em 2012},~930849.
\newblock {\url{https://doi.org/10.1100/2012/930849}}.

\bibitem[{Analytica Laboratories}(2020)]{analytica2020}
{Analytica Laboratories}.
\newblock Managing the Risk of Fermenting Honey.
\newblock Technical Report; Analytica Laboratories.  2020.  Available online:  \url{https://www.analytica.co.nz/media/nuwbyumm/managing-the-risk-of-furmenting-honey.pdf}  \hl{(accessed on 31 May 2026.).} %MDPI: Please add the access date (format: Date Month Year), e.g., accessed on 1 January 2020. Please notice that the accessed date should be before the accepted date.
%Date added


\bibitem[Gleiter et~al.(2006)Gleiter, Horn, and Isengard]{gleiter2006}
Gleiter, R.A.; Horn, H.; Isengard, H.D.
\newblock {Influence of Type and State of Crystallisation on the Water Activity
  of Honey}.
\newblock {\em Food Chem.} {\bf 2006}, {\em 96},~441--445.
\newblock {\url{https://doi.org/10.1016/j.foodchem.2005.02.051}}.

\bibitem[Magar et~al.(2019)]{ijfans2019}
Shafiq, H.; Iftikhar, F.; Ahmad, A.; Kaleem, M.; Sair, A.T.
\newblock {Effect of Crystallization on the Water Activity of Honey}.
\newblock \emph{Int. J. Food Nutr. Sci.}  \hl{\textbf{2014}, \emph{3}, 1--6.} 

%I changed the year from 2019 to 2014. The full details are:
%Shafiq, H.; Iftikhar, F.; Ahmad, A.; Kaleem, M.; Sair, A.T. Effect of Crystallization on the Water Activity of Honey. Int. J. Food %Nutr. Sci. 2014, 3 (3), 1–6.
%https://www.ijfans.org/uploads/paper/9edaa46c917db303412fd7200b88b447.pdf

\bibitem[{Apiculture New Zealand}(2021)]{apinz2021}
{Apiculture New Zealand}.
\newblock {\emph{Managing the Risk of Fermenting Honey}};
\newblock Apiculture NZ Industry Document: \hl{Wellington, New Zealand,} 
  2021.
%Thank you for the comment. We have identified that this reference (apinz2021) is a duplicate of another reference in our list (analytica2020) — both point to the same document, "Managing the Risk of Fermenting Honey" (New Zealand Beekeeper, November 2018). Please remove this duplicate (apinz2021) and rearrange the citation order accordingly. We confirm the remaining reference is cited in the main

\bibitem[Phiri et~al.(2022)]{phiri2022}
Phiri, B.J.; Fèvre, D.; Hidano, A.
\newblock {Uptrend in global managed honey bee colonies and production based on
  a six-decade viewpoint}.
\newblock {\em Sci. Rep.} {\bf 2022}, {\em 12},~2045--2322.
\newblock {\url{https://doi.org/10.1038/s41598-022-25290-3}}.

\bibitem[{Destatis---Statistisches Bundesamt}(2026)]{destatis2026}
{Destatis---Statistisches Bundesamt}.
\newblock {\emph{International Statistics: Global Honey Bee Population on the Rise}}; 
\newblock Destatis, Federal Statistical Office Germany: \hl{Wiesbaden,  Germany,} 
 2026.
 %Berling is incorrect. Changed to: Wiesbaden, Germany

\bibitem[Lastriyanto et~al.(2020)Lastriyanto, Wibowo, Erwan, Jaya, Batoro,
  Masyithoh, and Lamerkabel]{lastriyanto2020}
Lastriyanto, A.; Wibowo, S.A.; Erwan.; Jaya, F.; Batoro, J.; Masyithoh, D.;
  Lamerkabel, J.S.A.
\newblock Moisture Reduction of Honey in Dehumidification and Evaporation
  Processes.
\newblock {\em J. Mech. Eng. Sci. Technol.} {\bf
  2020}, {\em 4},~153--163.
\newblock {\url{https://doi.org/10.17977/um016v4i22020p153}}.

\bibitem[Ikhsan et~al.(2022)Ikhsan, Chin, and Ahmad]{spel3upm}
Ikhsan, L.N.; Chin, K.Y.; Ahmad, F.
\newblock Methods of the Dehydration Process and Its Effect on the
  Physicochemical Properties of Stingless Bee Honey: {A} Review. \emph{Molecules}   \textbf{2022}, \emph{27}, 7243.
\newblock {\url{https://doi.org/10.3390/molecules27217243}}.

\bibitem[Beyene et~al.(2015)Beyene, Abi, Chalchissa, and
  WoldaTsadik]{ethaau2007}
Beyene, T.; Abi, D.; Chalchissa, G.; WoldaTsadik, M.
\newblock Evaluation of Transitional and Modern Hives for Honey Production in
  Mid Rift Valley of Ethiopia. \emph{Glob. J. Anim. Sci. Res.}   \textbf{2015}, \emph{3}, 48--56.

\bibitem[Gobessa et~al.(2012)]{sciendo2012}
Gobessa, S.;  Seifu, E.;  Bezabih, A.
\newblock {Physicochemical Properties of Honey Produced in the Homesha District
  of Western Ethiopia}.
\newblock {\em J. Apic. Sci.} {\bf 2012}, \emph{56}, 56.
\newblock {\url{https://doi.org/10.2478/v10289-012-0004-z}}.

\bibitem[Echeverrigaray et~al.(2021)Echeverrigaray, Scariot, Foresti, Schwarz,
  Rocha, da~Silva, Moreira, and Delamare]{stinglessyeastbrazil2021}
Echeverrigaray, S.; Scariot, F.J.; Foresti, L.; Schwarz, L.V.; Rocha, R.K.M.;
  da~Silva, G.P.; Moreira, J.P.; Delamare, A.P.L.
\newblock Yeast biodiversity in honey produced by stingless bees raised in the
  highlands of southern {Brazil}.
\newblock {\em Int. J. Food Microbiol.} {\bf 2021}, {\em
  347},~109200.
\newblock {\url{https://doi.org/10.1016/j.ijfoodmicro.2021.109200}}.

\bibitem[Melia et~al.(2024)Melia, Juliyarsi, Kurnia, Aritonang, Rusdimansyah,
  Sukma, Setiawan, Pratama, and Supandil]{stinglesswjava2024}
Melia, S.; Juliyarsi, I.; Kurnia, Y.F.; Aritonang, S.N.; Rusdimansyah, R.;
  Sukma, A.; Setiawan, R.D.; Pratama, Y.E.; Supandil, D.
\newblock Profile of stingless bee honey and microbiota produced in {West
  Sumatra}, {Indonesia}, by several species ({Apidae, Meliponinae}).
\newblock {\em Vet. World} {\bf 2024}, {\em 17},~785--795.
\newblock {\url{https://doi.org/10.14202/vetworld.2024.785-795}}.

\bibitem[Haftom et~al.(2013)]{lrrd2013}
Gebremedhn, H.;  Estifanos, A..
\newblock {On Farm Evaluation of Kenyan Top Bar Hive (KTBH) for Honey
  Production in Tigray, Ethiopia}.
\newblock {\em Livest. Res. Rural. Dev.} {\bf 2013}, {\em 25}, \hl{page.} %MDPI: please add page or doi information
%Thank you for the comment. Livestock Research for Rural Development uses article numbers rather than page numbers. The correct reference is Article #86, Volume 25 (5). Full details:
%Gebremedhn, H.; Estifanos, A. On Farm Evaluation of Kenyan Top Bar Hive (KTBH) for Honey Production in Tigray, Ethiopia. Livest. Res. Rural Dev. 2013, 25 (5), Article 86. Available online: https://www.lrrd.org/lrrd25/5/haft25086.htm

\bibitem[{European Commission, DG Agriculture and Rural
  Development}(2024)]{ec_honey_autumn2024}
{European Commission; DG Agriculture and Rural Development}.
\newblock \emph{Honey Market Presentation}; 
\newblock Technical report; European Commission: \hl{Brussels, Belgium,} 
  2024.
%Thank you for the comment. We apologize - the location should be Brussels, Belgium, not Ispra, Italy. DG Agriculture and Rural Development is based in Brussels. Corrected reference:
%European Commission, DG Agriculture and Rural Development. Honey Market Presentation; Technical report; European Commission: Brussels, Belgium, 2024.

\bibitem[{Eurostat}(2024)]{eurostat_honey2023}
{Eurostat}.
\newblock {EU} Imported \texteuro 359.3 Million Worth of Honey in 2023.
\newblock Eurostat Daily Data News, World Bee Day Release.  2024.  Available online:  \url{https://ec.europa.eu/eurostat/web/products-eurostat-news/w/ddn-20240520-1}  \hl{(accessed on 31 May 2026).} %MDPI: Please add the access date (format: Date Month Year), e.g., accessed on 1 January 2020. Please notice that the accessed date should be before the accepted date.
%Date added


\bibitem[{Tridge Global Intelligence}(2024)]{tridge2024}
{Tridge Global Intelligence}.
\newblock {\emph{China Is the Largest Supplier of Honey to the EU: 37\% of Extra-EU
  Imports in 2023}}; 
\newblock Tridge Market Intelligence: \hl{Seoul, Republic of Korea,} 
  2024.
%I confirm

\bibitem[{European Commission, DG SANTE}(2022)]{eufromhives2021}
{European Commission, DG SANTE}.
\newblock {\emph{EU Coordinated Action `From the Hives' (Honey 2021--2022): 46\% of
  Imported Honey Suspicious of Adulteration}};
\newblock European Commission, Food Safety: \hl{Brussels, Belgium,} 
  2022.
%Thank you for the comment. We apologize — the location should be Brussels, Belgium, not Ispra, Italy. DG SANTE (Health and Food Safety) is based in Brussels. Corrected reference:
%European Commission, DG SANTE. EU Coordinated Action 'From the Hives' (Honey 2021–2022): 46% of Imported Honey Suspicious of Adulteration; European Commission, Food Safety: Brussels, Belgium, 2022.

\bibitem[{European Parliament and Council of the European
  Union}(2024)]{cbi_traceability2024}
{European Parliament and Council of the European Union}.
\newblock Directive ({EU}) 2024/1438 of the {European Parliament} and of the
  Council of 14~{May} 2024 amending {Council} {Directives} 2001/110/{EC}
  relating to honey, 2001/112/{EC} relating to fruit juices, 2001/113/{EC}
  relating to fruit jams, jellies and marmalades, and 2001/114/{EC} relating to
  certain partly or wholly dehydrated preserved milk for human consumption.
\newblock \emph{Off. J. Eur. Union}  \textbf{2024}, \emph{L}, 1438.

\bibitem[{European Commission}(2023)]{EuropeanCommission2023NMR}
{European Commission}.
\newblock Response to Written Question {E-001756/2023} on the Use of {1H-NMR}
  Profiling in Honey Authenticity Control (Reply by Commissioner Stella
  Kyriakides on Behalf of the {European Commission}).
\newblock European Parliament---Written questions.  2023.  Available online: \url{https://www.europarl.europa.eu/RegData/seance_pleniere/textes_deposes/interpellations/2023/01756/P9_I(2023)01756_EN.pdf}  (accessed on 7 May 2026).


\bibitem[{\'A}lvarez-Su{\'a}rez et~al.(2025){\'A}lvarez-Su{\'a}rez, Majtan,
  Tejera, Santos-Buelga, and Gonz{\'a}lez-Param{\'a}s]{honeyid2025}
{\'A}lvarez-Su{\'a}rez, J.M.; Majtan, J.; Tejera, E.; Santos-Buelga, C.;
  Gonz{\'a}lez-Param{\'a}s, A.M.
\newblock The molecular identity of honey: Toward reliable biochemical
  authentication.
\newblock {\em Trends Food Sci. Technol.} {\bf 2025}, {\em
  155},~105331.
\newblock {\url{https://doi.org/10.1016/j.tifs.2025.105331}}.

\end{thebibliography}




%%%%%%%%%%%%%%%%%%%%%%%%%%%%%%%%%%%%%%%%%%
\PublishersNote{}
\end{adjustwidth}
\end{document}

